\documentclass[11pt]{article}
\usepackage{tcee}
\addbibresource{3A32.bib}
\tcetema[Oferta y demanda agregadas en una economía abierta]{3.A.32}{El modelo de oferta y demanda agregada: determinación de renta e inflación en una economía abierta. Análisis de las políticas monetaria y fiscal, y de los shocks y políticas de oferta.}
\tcefecha{13/11/2026}
\begin{document}
\tceetitulo

\section*{Introducción}
\minutos{2,5}

\begin{itemize}
\item \textbf{\emph{Enganche}}:
  \begin{itemize}
  \item «Las desgracias nunca vienen solas», dice el refrán.
  \item En 2022 a la zona del euro le llegaron 2 a la vez. Con la invasión de Ucrania se disparó el precio de la energía que importa, y el euro se depreció frente al dólar, la moneda en que se paga esa energía. Entre enero de 2021 y junio de 2022 el barril de Brent subió un 124~\% en dólares, pero un 158~\% en euros\nocite{EIA2026Brent,BCE2026EXR}.
  \item Y llegaron sobre una economía que salía de la Covid, con la demanda recuperándose deprisa y las cadenas de suministro atascadas. La inflación de la zona del euro llegó al 10,6~\% en octubre de 2022\nocite{Eurostat2026IAPC}. En 2026 la historia se repite: con el conflicto de Oriente Medio, el Brent llegó a 123,9 dólares en octubre, por encima del máximo de 2022, y la inflación volvió al 4,0~\%\nocite{EIA2026Brent,Eurostat2026IAPCoct}.
  \item Detrás están las preguntas de esta exposición: \emph{¿qué le pasa a la renta y a la inflación de una economía abierta cuando la golpea una perturbación, y qué deben hacer entonces la política monetaria, la fiscal y las políticas de oferta?}
  \end{itemize}

\item \textbf{\emph{Relevancia}}:
  \begin{itemize}
  \item \emph{Para la teoría}: el modelo de oferta y demanda agregadas es el marco con el que se piensan el corto y el medio plazo. Su versión moderna, el modelo neokeynesiano, es el que usan hoy los bancos centrales, y en una economía abierta añade 2 precios, el tipo de cambio y el tipo de interés mundial.
  \item \emph{Para la política económica}: la respuesta correcta depende del \emph{origen} de la perturbación (demanda u oferta) y, en una economía abierta, del régimen cambiario.
    \begin{itemize}
    \item España es una economía abierta: en 2025 sus exportaciones de bienes y servicios fueron el 36,5~\% del PIB y sus importaciones el 32,6~\%, una apertura del 69~\%\nocite{INE2026CNA2025,Eurostat2026nama10gdp}.
    \item España no tiene tipo de cambio propio frente a sus socios del euro, pero la zona del euro flota frente al dólar: los 2 regímenes conviven en la misma economía.
    \end{itemize}
  \item \emph{Para la actualidad}: desde 2020 se han sucedido una perturbación de oferta que acabó siendo de demanda (la Covid), una de demanda que chocó con una oferta atascada (2021) y otra de oferta externa (la energía de 2022 y de 2026), como en los 70. Las políticas de oferta están en el centro de la agenda europea\nocite{CE2024Draghi}.
  \end{itemize}

\item \textbf{\emph{Contextualización}}:
  \begin{itemize}
  \item \emph{Las definiciones de partida}:
    \begin{itemize}
    \item una \concepto{perturbación} (\emph{shock}) es un cambio exógeno e imprevisto. Es \emph{de demanda} si mueve en el mismo sentido la producción y los precios, y \emph{de oferta} si los mueve en sentido contrario. Se \emph{propaga} cuando los agentes reaccionan y la transmiten a otros mercados y periodos, y tiene \emph{efectos de segunda ronda} cuando contagia a los salarios y a los márgenes. Una \concepto{política de oferta} es la que aumenta la capacidad productiva o reduce los costes;
    \item una \concepto{economía abierta} comercia bienes y servicios con el resto del mundo y presta y se endeuda con él; su contabilidad la recoge la balanza de pagos \vertema{3.B.11}. Es \concepto{pequeña} si no influye en la renta ni en los tipos de interés del resto del mundo;
    \item el \concepto{tipo de cambio nominal}, $E$, es el precio de la moneda extranjera en moneda nacional (euros por dólar): si sube, la moneda nacional se \emph{deprecia}. El \concepto{tipo de cambio real}, $q\equiv EP^*/P$, compara los precios de fuera con los de dentro: si sube, se gana competitividad \vertema{3.B.13};
    \item el \concepto{régimen cambiario} dice quién fija el tipo de cambio: con \emph{tipo fijo}, el banco central se compromete a comprar y vender divisas a un precio dado; con \emph{tipo flexible}, lo fija el mercado \vertema{3.B.15}. La \concepto{movilidad del capital} mide cuánto responden los flujos financieros al diferencial de tipos de interés.
    \end{itemize}
  \item \emph{De dónde viene el modelo}:
    \begin{itemize}
    \item El IS-LM de \autor{John Hicks}\footnote{\autor{Hicks} recibió el Premio Nobel de Economía en 1972, con \autor{Kenneth Arrow}, «por sus contribuciones pioneras a la teoría del equilibrio económico general y a la teoría del bienestar»\nocite{Nobel1972HicksArrow}.} \vertema{3.A.5} lo llevaron a la economía abierta \textcite{Fleming1962Domestic} y \textcite{Mundell1963Capital}.\footnote{\autor{Robert Mundell} recibió el Premio Nobel de Economía en 1999 «por su análisis de la política monetaria y fiscal con distintos regímenes cambiarios y su análisis de las áreas monetarias óptimas»\nocite{Nobel1999Mundell}.}
    \item En los 70, con las crisis del petróleo, los precios dejaron de poder darse por fijos: la estanflación obligó a completar el modelo con la oferta agregada y con las perturbaciones de oferta.
    \item Desde los 90, el modelo neokeynesiano lo ha reconstruido con agentes que optimizan, expectativas racionales y precios rígidos\nocite{ClaridaGaliGertler1999Science,Woodford2003Interest}, también para la economía abierta\nocite{ObstfeldRogoff1995Redux,GaliMonacelli2005Monetary}. Es lo que premió el Nobel de 2026.\footnote{\autor{Michael Woodford} recibió el Premio Nobel de Economía en 2026, con \autor{Robert E. Hall}, «por los fundamentos microeconómicos de la estabilización macroeconómica»\nocite{Nobel2026Prensa}.}
    \end{itemize}
  \end{itemize}

\item \textbf{\emph{Problemática}}:
  \begin{itemize}
  \item ¿Cómo se determinan la renta y la inflación en una economía abierta, y qué modelos sirven para analizarlo?
  \item ¿Cómo se originan y se propagan las perturbaciones de oferta y de demanda, y cómo debe responder la política económica a cada una?
  \item ¿Qué cambia cuando la perturbación llega de fuera y la economía tiene que elegir régimen cambiario?
  \end{itemize}

\item \textbf{\emph{Estructura}}:
  \begin{itemize}
  \item La exposición sigue la vida de una perturbación: primero el marco en el que cae, después de dónde viene, cómo se propaga y cómo se responde, y al final qué cambia cuando viene de fuera (figura~\ref{figu:estructura-tema}):
    \begin{itemize}
    \item \concepto{I. El marco} (\ref{sec:marco}): del modelo de \autor{Mundell}-\autor{Fleming}, con precios fijos, a la demanda y la oferta agregadas abiertas, que dejan mover los precios, y de ahí a su versión con microfundamentos, el modelo neokeynesiano. Cada modelo responde a un límite del anterior, y en cada uno se ve qué pueden la política monetaria y la fiscal;
    \item \concepto{II. Las perturbaciones de oferta y de demanda} (\ref{sec:perturbaciones}): con ese marco, el origen, la propagación y la respuesta de política para cada tipo, primero las de oferta y después las de demanda, que a menudo nacen de ellas;
    \item \concepto{III. Extensión: las perturbaciones en una economía abierta} (\ref{sec:abierta}): el modelo neokeynesiano abierto y la respuesta a las perturbaciones que llegan de fuera según el régimen cambiario.
    \end{itemize}
  \end{itemize}
\end{itemize}

\begin{esquema}[unbreakable]
\textbf{3.A.32. Oferta y demanda agregadas en una economía abierta}\\[2pt]
Introducción\esqtiempo{2,5}\\
I. El marco: de \autor{Mundell}-\autor{Fleming} al modelo neokeynesiano\esqtiempo{11,5}\\
\hspace*{1.5em}1. El modelo de \autor{Mundell}-\autor{Fleming}\\
\hspace*{1.5em}2. La demanda y la oferta agregadas en una economía abierta\\
\hspace*{1.5em}3. El modelo neokeynesiano\\
II. Las perturbaciones de oferta y de demanda\esqtiempo{10}\\
\hspace*{1.5em}1. Las de oferta: origen y propagación\\
\hspace*{1.5em}2. La respuesta a las de oferta\\
\hspace*{1.5em}3. Las de demanda y su interacción con las de oferta\\
\hspace*{1.5em}4. La respuesta a las de demanda\\
III. Extensión: las perturbaciones en una economía abierta\esqtiempo{5}\\
\hspace*{1.5em}1. El modelo neokeynesiano de una economía abierta\\
\hspace*{1.5em}2. Las perturbaciones externas y el régimen cambiario\\
Conclusión\esqtiempo{1}
\end{esquema}

\begin{notaopositor}
Reparto del cante: introducción, \tiempo{2,5}; I, \tiempo{11,5} (I.1, \tiempo{5}; I.2, \tiempo{3,5}; I.3, \tiempo{3}); II, \tiempo{10} (II.1, \tiempo{2,5}; II.2, \tiempo{3}; II.3, \tiempo{2}; II.4, \tiempo{2,5}); III, \tiempo{5} (III.1, \tiempo{2}; III.2, \tiempo{3}), y conclusión, \tiempo{1}. Todo el tema vuelve a una misma economía: España, pequeña y abierta, con el tipo de cambio fijo frente a sus socios del euro y flotante frente al dólar, ante la Covid y la energía cara de 2022 y de 2026. La movilidad del capital se escribe $\beta_K$, porque en el modelo neokeynesiano $\beta$ es el factor de descuento. Las derivaciones están en los anexos~\ref{anexo:mundell-fleming} a~\ref{anexo:nek}; el modelo de redes y el cuadro de perturbaciones y respuestas, en los anexos~\ref{anexo:redes} y~\ref{anexo:perturbaciones}.
\end{notaopositor}

\begin{infografia}[p]
\centering
\caption{La estructura del tema: el marco, las perturbaciones y la economía abierta}\label{figu:estructura-tema}
\grafico[\textwidth]{estructura-tema}
\fuente{Elaboración propia. Cada modelo del bloque~\ref{sec:marco} responde a un límite del anterior; el bloque~\ref{sec:perturbaciones} usa ese marco para cada tipo de perturbación (origen y propagación, y después respuesta de política), y el bloque~\ref{sec:abierta} lo extiende a las perturbaciones que llegan de fuera.}
\end{infografia}


% ==========================================================================
\section{El marco: de \autor{Mundell}-\autor{Fleming} al modelo neokeynesiano}\label{sec:marco}
\minutos{11,5}

\begin{itemize}
\item Empezamos por el marco en el que caen las perturbaciones: cómo se determinan la renta y la inflación en una economía abierta y qué pueden hacer con ellas la política monetaria y la fiscal.
\item \emph{Conclusión anticipada}: 3 modelos, cada uno para resolver un límite del anterior.
  \begin{itemize}
  \item Con precios fijos, el modelo de \autor{Mundell} y \autor{Fleming} determina la renta y muestra que la eficacia de cada política depende del régimen cambiario y de la movilidad del capital.
  \item Si se deja mover el nivel de precios, cada equilibrio de \autor{Mundell}-\autor{Fleming} da un punto de la demanda agregada, y con la oferta agregada se determinan la renta y los precios; la economía abierta añade la competitividad a la demanda y la cuña exterior a la oferta.
  \item Con microfundamentos y expectativas, la demanda y la oferta agregadas se convierten en la IS dinámica y la curva de \autor{Phillips} neokeynesiana, que separan con claridad las perturbaciones de demanda de las de costes.
  \end{itemize}
\end{itemize}

% --------------------------------------------------------------------------
\subsection{El modelo de \autor{Mundell}-\autor{Fleming}}\label{sec:mundell-fleming}

\paragraph{a) Idea}

\begin{itemize}
\item A principios de los 60, con el sistema de Bretton Woods en marcha y los capitales empezando a moverse, \textcite{Fleming1962Domestic}, desde el FMI, y \textcite{Mundell1963Capital} se preguntaron qué eficacia tienen la política monetaria y la fiscal cuando el capital es móvil, con tipo de cambio fijo y con tipo flexible.
\item Su respuesta fue el modelo IS-LM-BP: el IS-LM \vertema{3.A.5} con el comercio en la curva IS y un tercer mercado, el de divisas, resumido en la curva BP.
\item Fue la referencia de la macroeconomía abierta durante 30 años y sigue siendo la intuición de partida de los bancos centrales y del FMI\nocite{Blanchard2021Macroeconomics}. Aquí es, además, el primer peldaño: la demanda agregada abierta (\ref{sec:da-oa}) se construye con él.
\end{itemize}

\paragraph{b) Modelo}

\subparagraph{i) Supuestos}

\begin{itemize}
\item \emph{Precios fijos}: $P=P^*=1$. No hay inflación, el tipo nominal es el real ($i=r$) y el tipo de cambio real es el nominal ($q=E$).
\item \emph{País pequeño}: la renta exterior, $Y^*$, y el tipo de interés mundial, $i^*$, están dados.
\item \emph{Expectativas estáticas} del tipo de cambio: no se espera que cambie ($\gamma E^e=0$), de modo que la paridad no cubierta de intereses, $i=i^*+\gamma E^e$, se reduce a $i=i^*$ cuando el capital es perfectamente móvil \vertema{3.B.13}.
\item \emph{Activos nacionales y extranjeros sustitutos perfectos} y agentes neutrales al riesgo: no hay prima de riesgo.
\item \emph{Condición de \autor{Marshall}-\autor{Lerner}}: una depreciación real mejora el saldo comercial \vertema{3.B.12}.\footnote{A corto plazo puede empeorarlo antes de mejorarlo, la \emph{curva J}: los pedidos se cierran con meses de antelación, de modo que los volúmenes responden todavía al tipo de cambio antiguo mientras la depreciación ya encarece, en moneda nacional, lo importado. En los países industriales suele durar entre 6 meses y un año\nocite{KrugmanObstfeldMelitz2018International}.}
\item \emph{Los flujos de capital responden al diferencial de tipos}: la cuenta financiera neta es $\beta_K\,(i-i^*)$, con $\beta_K\geq0$ la \concepto{movilidad del capital}.\footnote{En esta exposición $\beta_K$ lleva subíndice porque $\beta$ es el factor de descuento del modelo neokeynesiano (\ref{sec:nek}). Es un enfoque de \emph{flujos}: mientras haya diferencial, entra capital. Su crítica está en la valoración del modelo.} No es un supuesto sobre su valor: el modelo se resuelve para los 3 casos.
\item \emph{Sin esterilización}: con tipo fijo, lo que el banco central compra o vende de divisas cambia la base monetaria.
\end{itemize}

\subparagraph{ii) Desarrollo}

\begin{itemize}
\item \emph{Mercado de bienes: la curva IS abierta}. A la demanda interna se suman las exportaciones netas, $XN=x\,Y^*-m\,Y+\nu\,q$: crecen con la renta exterior ($x$, propensión a exportar), caen con la propia ($m$, propensión a importar) y crecen con el tipo de cambio real ($\nu$, sensibilidad del saldo comercial). Con $\overline{A}\equiv\overline{C}+\overline{I}+\overline{G}$ el gasto autónomo,
  \funcionclave[eq:is-abierta]{\text{IS:}\quad Y=\frac{\overline{A}-b\,i+x\,Y^*+\nu\,q}{1-c\,(1-t)+m}}
  \begin{itemize}
  \item \emph{En palabras}: es la IS del modelo cerrado con 2 cambios. El multiplicador es menor, porque parte del gasto se va en importaciones ($m$ en el denominador), y la curva se desplaza con la renta exterior y con el tipo de cambio: una depreciación ($q\uparrow$) la lleva a la derecha.
  \item Volvamos a España: su producción depende de lo que crecen sus socios ($Y^*$), del tipo de cambio real frente a ellos ($q$) y del tipo de interés que se fija fuera ($i^*$). Con una apertura del 69~\%, una caída de la demanda de la zona del euro se nota enseguida en su renta\nocite{INE2026CNA2025}.
  \end{itemize}
\item \emph{Mercado de dinero: la curva LM}, la de la economía cerrada:
  \begin{equation}
  \text{LM:}\quad \frac{M}{P}=k\,Y-h\,i . \label{eq:lm}
  \end{equation}
  Con tipo fijo, la cantidad de dinero ya no la decide el banco central: es el crédito interno más las reservas, $M=\mathrm{CI}+R$, y cambia con cada intervención en el mercado de divisas.
\item \emph{Mercado de divisas: la curva BP}. Son los pares $(Y,i)$ con los que la balanza de pagos está en equilibrio, sin variación de reservas: el saldo comercial se financia con los flujos de capital,
  \funcionclave[eq:bp]{\text{BP:}\quad \underbrace{x\,Y^*-m\,Y+\nu\,q}_{\text{saldo comercial}}+\underbrace{\beta_K\,(i-i^*)}_{\text{cuenta financiera neta}}=0\quad\Longrightarrow\quad \frac{\mathrm{d}i}{\mathrm{d}Y}\bigg|_{\mathrm{BP}}=\frac{m}{\beta_K}}
  \begin{itemize}
  \item \emph{En palabras}: con más renta se importa más y el saldo comercial empeora; para financiarlo hace falta un tipo de interés más alto que atraiga capital. Por encima de la BP hay superávit (entran divisas); por debajo, déficit. Una depreciación la desplaza a la derecha.
  \item \emph{Los 3 casos} (gráfico~\ref{fig:is-lm-bp}):
    \begin{itemize}
    \item \emph{movilidad nula} ($\beta_K=0$): la BP es vertical; solo hay una renta compatible con el equilibrio comercial;
    \item \emph{movilidad imperfecta} ($0<\beta_K<\infty$): la BP es creciente; la movilidad es \emph{alta} si la BP es más plana que la LM ($m/\beta_K<k/h$) y \emph{baja} si es más inclinada;
    \item \emph{movilidad perfecta} ($\beta_K\to\infty$): la BP es horizontal en $i=i^*$; el menor diferencial atraería o expulsaría capital sin límite.
    \end{itemize}
  \end{itemize}
\end{itemize}

\begin{figure}[H]
\centering
\caption{El modelo IS-LM-BP: equilibrio y los 3 casos de movilidad del capital}\label{fig:is-lm-bp}
\grafico[0.62\textwidth]{is-lm-bp}
\fuente{Elaboración propia a partir de \textcite{Fleming1962Domestic}, \textcite{Mundell1963Capital} y \textcite{Gandolfo2002International}. Por encima de cada BP, superávit de la balanza de pagos; por debajo, déficit.}
\end{figure}

\begin{itemize}
\item \emph{Equilibrio}: las 3 curvas se cortan en un punto. Son 3 ecuaciones, y lo que se ajusta para que se corten depende del régimen:
  \begin{itemize}
  \item \emph{con tipo fijo}, las incógnitas son la renta, el tipo de interés y la \emph{cantidad de dinero}: si hay superávit, el banco central compra divisas para que la moneda no se aprecie, la base monetaria crece y la LM se desplaza hasta pasar por el corte de la IS y la BP;
  \item \emph{con tipo flexible}, las incógnitas son la renta, el tipo de interés y el \emph{tipo de cambio}: si hay superávit, la moneda se aprecia, las exportaciones netas caen y la IS y la BP se desplazan hasta cortarse sobre la LM, que controla el banco central.
  \end{itemize}
\item \emph{Las políticas con movilidad perfecta}. Con $\beta_K\to\infty$, el tipo de interés no puede separarse del mundial. Son 4 casos (gráfico~\ref{fig:mundell-fleming-politicas}):
  \begin{itemize}
  \item \emph{Tipo fijo, política fiscal: muy eficaz}. Más gasto desplaza la IS a la derecha y presiona al alza sobre el tipo de interés; entra capital y la moneda tiende a apreciarse. Para evitarlo, el banco central compra divisas, la cantidad de dinero crece y la LM acompaña a la IS. No hay expulsión por el tipo de interés: la renta sube con todo el multiplicador abierto, $1/[1-c\,(1-t)+m]$.
  \item \emph{Tipo fijo, política monetaria: ineficaz}. Más dinero baja el tipo de interés, sale capital y la moneda tiende a depreciarse; el banco central vende divisas, la cantidad de dinero vuelve a su nivel y la LM a su sitio. Solo cambia la composición del balance del banco central: más crédito interno y menos reservas.
  \item \emph{Tipo flexible, política fiscal: ineficaz}. La presión al alza sobre el tipo de interés atrae capital y aprecia la moneda; las exportaciones netas caen hasta que la IS vuelve a su sitio. La expulsión (\emph{crowding out}) no es de la inversión, sino de las exportaciones netas: caen exactamente lo que sube el gasto.
  \item \emph{Tipo flexible, política monetaria: muy eficaz}. Baja el tipo de interés, sale capital y la moneda se deprecia; las exportaciones netas suben y la IS se desplaza a la derecha hasta cortar a la LM nueva en $i=i^*$. La renta sube lo que pide el mercado de dinero, $\Delta Y=\Delta(M/P)/k$: la política actúa por el canal del tipo de cambio \vertema{3.A.37}.
  \end{itemize}
\item \emph{Una devaluación} con tipo fijo sube las exportaciones netas y desplaza la IS: funciona como una expansión fiscal, pero cambiando el gasto de bienes extranjeros a nacionales en vez de aumentarlo. Es lo que \textcite{Meade1951} llamó una política de \emph{cambio del gasto} (\emph{expenditure switching}), frente a las de \emph{nivel del gasto} (\emph{expenditure changing}).\footnote{\autor{James Meade} recibió el Premio Nobel de Economía en 1977, con \autor{Bertil Ohlin}, «por su contribución pionera a la teoría del comercio internacional y de los movimientos internacionales de capital»\nocite{Nobel1977MeadeOhlin}.}
\end{itemize}

\begin{figure}[H]
\centering
\caption{Políticas fiscal y monetaria expansivas con movilidad perfecta del capital}\label{fig:mundell-fleming-politicas}
\grafico[0.9\textwidth]{mundell-fleming-politicas}
\fuente{Elaboración propia a partir de \textcite{Fleming1962Domestic} y \textcite{Mundell1963Capital}. Arriba, tipo de cambio fijo; abajo, flexible. A la izquierda, expansión fiscal; a la derecha, monetaria. $A$, equilibrio inicial; $B$, el punto al que lleva la política antes de que reaccionen las reservas o el tipo de cambio; $C$, equilibrio final.}
\end{figure}

\begin{itemize}
\item \emph{Las políticas con movilidad imperfecta}. Con $0<\beta_K<\infty$ la conclusión se matiza, y lo que decide es si la BP es más o menos inclinada que la LM. Resolviendo las 3 ecuaciones (anexo~\ref{anexo:mundell-fleming}):
  \funcionclave[eq:multiplicadores-fijo]{\text{tipo fijo:}\quad \frac{\partial Y}{\partial\overline{G}}=\frac{\beta_K}{[1-c\,(1-t)]\,\beta_K+m\,(b+\beta_K)},\qquad \frac{\partial Y}{\partial M}=0}
  \funcionclave[eq:multiplicadores-flexible]{\text{tipo flexible:}\quad \frac{\partial Y}{\partial\overline{G}}=\frac{h}{D_X},\qquad \frac{\partial Y}{\partial(M/P)}=\frac{b+\beta_K}{D_X},\qquad D_X\equiv h\,[1-c\,(1-t)]+(b+\beta_K)\,k}
  \begin{itemize}
  \item \emph{En palabras}: con tipo fijo, la fiscal gana eficacia con la movilidad (de 0 sin movilidad al multiplicador abierto con movilidad perfecta) y la monetaria no sirve nunca, porque el dinero es endógeno. Con tipo flexible, la monetaria gana eficacia con la movilidad (del multiplicador de la economía cerrada a $1/k$) y la fiscal la pierde (del multiplicador cerrado a 0).
  \item \emph{Por qué con tipo flexible la economía parece cerrada}: el tipo de cambio mantiene en equilibrio la balanza de pagos, de modo que el saldo comercial es lo que sale o entra de capital, $XN=-\beta_K\,(i-i^*)$. La IS y la BP juntas son la IS de una economía cerrada con sensibilidad al tipo de interés $b+\beta_K$ (anexo~\ref{anexo:mundell-fleming}).
  \item La eficacia de cada política en los 8 casos está en la tabla~\ref{tab:mundell-fleming}.
  \end{itemize}
\end{itemize}

\begin{table}[H]
\centering
\caption{Eficacia de las políticas monetaria y fiscal en el modelo de \autor{Mundell}-\autor{Fleming}}\label{tab:mundell-fleming}
\small
\begin{tabular}{>{\raggedright\arraybackslash}m{1.9cm}>{\raggedright\arraybackslash}m{3.0cm}>{\centering\arraybackslash}m{2.1cm}>{\centering\arraybackslash}m{2.1cm}>{\centering\arraybackslash}m{2.1cm}>{\centering\arraybackslash}m{2.1cm}}
\toprule
\rowcolor{tceeazulpalido} \multicolumn{2}{l}{Movilidad del capital} & Perfecta & Alta & Baja & Nula\\
\midrule
\multirow{2}{=}{Fiscal} & Tipo de cambio fijo & \textcolor{tceeverde}{Muy eficaz} & \textcolor{tceeverde}{Eficaz} & \textcolor{tceedorado}{Poco eficaz} & \textcolor{tceerojo}{Ineficaz}\\
 & Tipo de cambio flexible & \textcolor{tceerojo}{Ineficaz} & \textcolor{tceedorado}{Poco eficaz} & \textcolor{tceeverde}{Eficaz} & \textcolor{tceeverde}{Muy eficaz}\\
\midrule
\multirow{2}{=}{Monetaria} & Tipo de cambio fijo & \textcolor{tceerojo}{Ineficaz} & \textcolor{tceerojo}{Ineficaz} & \textcolor{tceerojo}{Ineficaz} & \textcolor{tceerojo}{Ineficaz}\\
 & Tipo de cambio flexible & \textcolor{tceeverde}{Muy eficaz} & \textcolor{tceeverde}{Eficaz} & \textcolor{tceeverde}{Eficaz} & \textcolor{tceeverde}{Eficaz}\\
\bottomrule
\end{tabular}
\fuente{Elaboración propia a partir de \textcite{Heijdra2017Foundations} y de \eqref{eq:multiplicadores-fijo} y \eqref{eq:multiplicadores-flexible}. Movilidad alta: la BP es menos inclinada que la LM; baja: más inclinada. Con tipo fijo, la monetaria solo cambia las reservas, más deprisa cuanto mayor es la movilidad; con tipo flexible, su eficacia crece con la movilidad, y sin movilidad es la de la economía cerrada.}
\end{table}

\subparagraph{iii) Implicaciones}

\begin{itemize}
\item \emph{El régimen decide qué se ajusta}: con tipo fijo, el banco central pierde el control de la cantidad de dinero, que se vuelve endógena; con tipo flexible, lo recupera, y el tipo de cambio absorbe los desequilibrios.
\item \emph{Cada política en su régimen}: con capital móvil, la monetaria funciona con tipo flexible y la fiscal con tipo fijo.
\item \emph{El dilema del régimen}: con tipo fijo, la política monetaria no sirve con \emph{ninguna} movilidad del capital. Con movilidad nula, la expansión del crédito interno empeora el comercio y se va en reservas a medio plazo; con movilidad perfecta, en el acto. Si se insiste, se agotan las reservas: es la crisis cambiaria de primera generación de \textcite{Krugman1979Model}\footnote{\autor{Paul Krugman} recibió el Premio Nobel de Economía en 2008 «por su análisis de los patrones del comercio y de la localización de la actividad económica»\nocite{Nobel2008Krugman}.} \vertema{3.B.17}.
\item \emph{El trilema}: de los 3 objetivos (un tipo de cambio fijo, capitales que entran y salen sin trabas y un banco central que fija sus tipos para la economía propia), solo se pueden tener 2 a la vez; como dice el refrán, no se puede nadar y guardar la ropa\nocite{ObstfeldShambaughTaylor2005} \vertema{3.B.15}.
  \begin{itemize}
  \item Cada régimen es un lado del triángulo y renuncia al vértice opuesto (figura~\ref{figu:trilema}): la unión monetaria, a la autonomía; la flotación, al tipo fijo; los controles de capitales, a la movilidad.
  \item No obliga a elegir una esquina: se puede ceder la mitad de cada objetivo, como en una flotación gestionada, y lo óptimo suele ser una solución interior\nocite{Frankel1999NoSingle}.
  \item \textcite{PadoaSchioppa1988} añadió el libre comercio como 4.º objetivo, el \emph{cuarteto imposible} (\emph{inconsistent quartet}): fue el argumento para pasar del Sistema Monetario Europeo a la unión monetaria, que conserva el tipo fijo, la movilidad y el comercio y renuncia a la política monetaria propia.
  \end{itemize}
\end{itemize}

\begin{infografia}[!htb]
\centering
\caption{El trilema de la política monetaria en una economía abierta}\label{figu:trilema}
\grafico[0.45\textwidth]{trilema}
\fuente{Elaboración propia a partir de \textcite{ObstfeldShambaughTaylor2005}. En los vértices, los 3 objetivos; en cada lado, el régimen que se queda con 2 de ellos y renuncia al vértice opuesto.}
\end{infografia}

\begin{itemize}
\item \emph{Aislamiento}: con tipo flexible, una caída de la demanda exterior no mueve la renta, porque la depreciación compensa la pérdida de exportaciones (anexo~\ref{anexo:mundell-fleming}). Es el argumento clásico de \textcite{Friedman1953Case} a favor de la flotación.\footnote{\autor{Milton Friedman} recibió el Premio Nobel de Economía en 1976 «por sus logros en el análisis del consumo, la historia y la teoría monetarias y por demostrar la complejidad de la política de estabilización»\nocite{Nobel1976Friedman}.} No aísla, en cambio, de las perturbaciones financieras: una subida del tipo de interés mundial mueve la renta en los 2 regímenes (\ref{sec:externas}).
\item \emph{Un instrumento para cada objetivo}: \textcite{Meade1951} planteó el conflicto entre el \emph{equilibrio interno} (pleno empleo sin inflación) y el \emph{externo} (equilibrio de la balanza de pagos), y \textcite{Tinbergen1952Theory} mostró que hacen falta tantos instrumentos como objetivos.\footnote{\autor{Jan Tinbergen} recibió el Premio Nobel de Economía en 1969, con \autor{Ragnar Frisch}, «por desarrollar y aplicar modelos dinámicos para el análisis de los procesos económicos»\nocite{Nobel1969}.} \textcite{Swan1963Longer} lo dibujó con el gasto interno y el tipo de cambio real como instrumentos.
  \begin{itemize}
  \item \emph{La asignación de \autor{Mundell}}: con tipo fijo y sin tocar el tipo de cambio, \textcite{Mundell1962Appropriate} propuso asignar la política monetaria al equilibrio externo y la fiscal al interno, porque la monetaria mueve más la balanza de pagos (a través del capital) en relación con lo que mueve la renta. Con la asignación al revés, cada corrección aleja la economía de su objetivo: es el principio de la clasificación eficaz de los mercados \vertema{3.B.12}.
  \end{itemize}
\item \emph{Los déficits gemelos}: más déficit público empeora la cuenta corriente por la renta (más importaciones; domina con tipo fijo) y por el tipo de interés y el tipo de cambio (entrada de capitales, apreciación y menos exportaciones netas; con tipo flexible y movilidad perfecta, las exportaciones netas caen tanto como sube el gasto) \vertema{3.B.11}. Fue la combinación de los Estados Unidos de principios de los 80: expansión fiscal, política monetaria contractiva, dólar fuerte y un déficit exterior creciente; las exportaciones netas pasaron del $-0{,}6$~\% del PIB en 1982 al $-2{,}5$~\% en 1984\nocite{FRED2026ExportacionesNetas}.
\item Volvamos a España. Con la peseta tenía política monetaria propia y podía devaluar, como hizo en 1992 y 1993 dentro del Sistema Monetario Europeo. Dentro del euro, frente a sus socios su tipo de cambio es fijo e irrevocable, el capital se mueve libremente y la política monetaria la decide el BCE para el conjunto: el trilema resuelto renunciando a la política monetaria propia. Para estabilizar lo propio le queda la política fiscal, que con capital móvil es muy eficaz, salvo cuando la deuda alta encarece la financiación \vertema{3.A.40} \vertema{4.B.4}.
\end{itemize}

\begin{anotaciones}
\textbf{Los regímenes cambiarios en el mundo}
\begin{itemize}
\item Según la clasificación de hecho del FMI, a 30 de abril de 2023, el 44,8~\% de sus miembros tenía un anclaje blando (\emph{soft peg}), el 32,5~\% flotaba, el 13,4~\% tenía un anclaje duro (\emph{hard peg}) y el 9,3~\% otros regímenes gestionados\nocite{FMI2024AREAER2023}.
\item Los países del euro aparecen en libre flotación: el euro flota frente al resto del mundo, aunque entre ellos no haya tipo de cambio \vertema{3.B.15}.
\end{itemize}
\end{anotaciones}

\paragraph{c) Extensiones}

\begin{itemize}
\item \emph{2 países grandes}: si el país es grande, su política mueve el tipo de interés mundial y la renta de los demás\nocite{Mundell1964Reply,Heijdra2017Foundations} \vertema{3.A.37} \vertema{3.B.19}:
  \begin{itemize}
  \item con salarios nominales rígidos en los 2, una expansión fiscal sube el tipo mundial, aprecia la moneda propia y desvía demanda hacia el socio: la renta sube en los 2, es una política de \emph{locomotora};
  \item una expansión monetaria deprecia la moneda propia y atrae demanda a costa del socio: es una política de \emph{empobrecer al vecino} (\emph{beggar-thy-neighbour});
  \item con salarios \emph{reales} rígidos los papeles se invierten: la fiscal empobrece al vecino y la monetaria deja de tener efectos reales.
  \end{itemize}
\item \emph{Expectativas}: si los agentes anticipan el tipo de cambio de largo plazo y los precios se ajustan despacio, el tipo de cambio reacciona de más a corto plazo, la \emph{sobrerreacción} (\emph{overshooting})\nocite{Dornbusch1976Expectations} \vertema{3.B.14}.
\item \emph{Fondos frente a flujos}: el enfoque monetario de la balanza de pagos\nocite{FrenkelJohnson1976Monetary} y los modelos de cartera tratan la balanza de pagos como el ajuste de los \emph{saldos} de dinero y de activos, no como un flujo permanente \vertema{3.B.14}.
\end{itemize}

\paragraph{d) Evidencia empírica}

\begin{itemize}
\item \emph{El trilema se ve en los datos}: con más de un siglo de datos, los países con tipo fijo y libre movilidad del capital siguen los tipos de interés del país de referencia; los que tienen controles de capitales o flotan conservan más autonomía monetaria\nocite{ObstfeldShambaughTaylor2005}. Con índices de los 3 objetivos desde 1970, su suma ponderada es aproximadamente constante: lo que se gana en uno se pierde en los otros\nocite{AizenmanChinnIto2008}.
\item \emph{La política fiscal, según el régimen}: con tipo de cambio flexible el multiplicador del gasto es casi nulo y con tipo fijo, positivo; en las economías más abiertas, menor\nocite{IlzetzkiMendozaVegh2013} \vertema{3.A.40}.
\item \emph{La movilidad del capital no es perfecta}: el ahorro y la inversión nacionales están muy correlacionados entre países, lo que no debería ocurrir si el capital se moviera sin fricciones\nocite{FeldsteinHorioka1980} \vertema{3.B.11}.
\end{itemize}

\paragraph{e) Valoración}

\begin{itemize}
\item \emph{Lo que aporta}: con 3 curvas, explica por qué la misma política funciona en un país y no en otro, y es la base del trilema y de las uniones monetarias \vertema{3.B.15} \vertema{3.B.16}.
\item \emph{Sus límites}:
  \begin{itemize}
  \item \emph{precios fijos}: no dice nada de la inflación, que es la mitad del título;
  \item \emph{flujos y no fondos}: un diferencial de tipos atrae capital una vez, hasta recomponer las carteras, no para siempre;
  \item \emph{expectativas estáticas y sin microfundamentos}: no tiene restricción intertemporal, de modo que no distingue un déficit exterior sostenible de uno que no lo es\nocite{ObstfeldRogoff1995Redux}.
  \end{itemize}
\item El primer límite es el más urgente para el título: dejamos que los precios se muevan.
\end{itemize}

% --------------------------------------------------------------------------
\subsection{La demanda y la oferta agregadas en una economía abierta}\label{sec:da-oa}

\paragraph{a) Idea}

\begin{itemize}
\item El modelo que da título al tema nace de levantar el primer supuesto de \autor{Mundell}-\autor{Fleming}, el de los precios fijos.
\item Como en la economía cerrada \vertema{3.A.5}, basta con dejar libre el nivel de precios: a cada $P$ le corresponde un equilibrio IS-LM-BP y, con él, una producción demandada; el conjunto de esos puntos es la \concepto{demanda agregada} (DA). La \concepto{oferta agregada} (OA) sale del mercado de trabajo. La economía abierta cambia las 2: la demanda gana el canal de la competitividad y la oferta paga lo que se compra fuera.
\end{itemize}

\paragraph{b) Modelo}

\subparagraph{i) Supuestos}

\begin{itemize}
\item Los de \autor{Mundell}-\autor{Fleming} (país pequeño, \autor{Marshall}-\autor{Lerner}, sin prima de riesgo), con \emph{movilidad perfecta del capital}, $i=i^*$, salvo que se diga otra cosa.
\item \emph{Nivel de precios nacional, $P$, variable} y exterior, $P^*$, dado: el tipo de cambio real, $q=EP^*/P$, ya no está fijado aunque lo esté el nominal.
\item \emph{Salarios nominales fijados con los precios esperados} y revisados con retraso, como en el modelo WS-PS \vertema{3.A.28}, pero negociados con el índice de precios de \emph{consumo}, $P_C$, del que una parte $\alpha$ es importada.\footnote{Aquí $\alpha$ es el grado de apertura del consumo, no la participación del capital de la función \autor{Cobb}-\autor{Douglas}. Es la misma $\alpha$ en todo el tema.}
\item \emph{Las empresas fijan el precio de lo que producen}, $P$, con un margen $\mu$ sobre sus costes, entre ellos los insumos importados (la energía).
\end{itemize}

\subparagraph{ii) Desarrollo}

\begin{itemize}
\item \emph{De la IS-LM-BP a la demanda agregada} (gráfico~\ref{fig:derivacion-da}). Un nivel de precios más alto tiene 2 efectos:
  \begin{itemize}
  \item el \emph{efecto saldos reales} o \emph{efecto \autor{Keynes}}: con más $P$, menos $M/P$, y la LM se desplaza a la izquierda, como en la economía cerrada;
  \item el \emph{efecto competitividad}, nuevo: con más $P$, menos $q$; los bienes nacionales se encarecen frente a los extranjeros, caen las exportaciones netas y la IS se desplaza a la izquierda.
  \end{itemize}
  Con movilidad perfecta, cada régimen se queda con uno de los 2 (anexo~\ref{anexo:oa-da}):
  \begin{itemize}
  \item \emph{con tipo fijo}, la LM sigue a la IS, porque el dinero es endógeno: manda el efecto competitividad (panel a);
  \item \emph{con tipo flexible}, la apreciación lleva a la IS hasta la nueva LM: manda el efecto saldos reales (panel b).
  \end{itemize}
  La demanda agregada de cada régimen es
  \funcionclave[eq:da-abierta]{\text{tipo fijo:}\ \ Y=\frac{\overline{A}-b\,i^*+x\,Y^*+\nu\,\overline{E}P^*/P}{1-c\,(1-t)+m};\qquad
  \text{tipo flexible:}\ \ Y=\frac{M/P+h\,i^*}{k}}
  \begin{itemize}
  \item \emph{En palabras}: con tipo fijo, la DA sale de la IS y se desplaza con la política fiscal, la demanda exterior y las devaluaciones, no con la cantidad de dinero; con tipo flexible, la DA sale de la LM, se desplaza con la política monetaria y no con la fiscal. Son los resultados de \autor{Mundell}-\autor{Fleming} pasados al plano $(Y,P)$.
  \item \emph{Sin movilidad y con tipo flexible}, el tipo de cambio mantiene equilibrado el comercio y la DA es la de la economía cerrada \vertema{3.A.5}.
  \end{itemize}
\end{itemize}

\begin{figure}[H]
\centering
\caption{De la IS-LM-BP a la demanda agregada abierta, con movilidad perfecta del capital}\label{fig:derivacion-da}
\grafico[0.85\textwidth]{derivacion-da-abierta}
\fuente{Elaboración propia a partir de \textcite{Blanchard2021Macroeconomics} y \textcite{Heijdra2017Foundations}. Arriba, el modelo IS-LM-BP; abajo, la DA. Al subir el nivel de precios de $P_0$ a $P_1$, el equilibrio pasa de $A$ a $B$. Con tipo fijo (a), el tipo real se aprecia, la IS se desplaza a la izquierda y la LM la sigue porque salen reservas; con tipo flexible (b), los saldos reales caen, la LM se desplaza a la izquierda y la apreciación lleva la IS al nuevo corte.}
\end{figure}

\begin{itemize}
\item \emph{La oferta agregada abierta: la cuña exterior}. A corto plazo, la OA es creciente, porque los salarios nominales tardan en revisarse \vertema{3.A.5}. La economía abierta añade 2 vías por las que el exterior entra en los costes:
  \begin{itemize}
  \item \emph{los insumos importados}: una subida del petróleo encarece producir; en el modelo de fijación de precios actúa como una subida del margen sobre el coste del trabajo y desplaza la OA hacia arriba\nocite{Blanchard2021Macroeconomics};
  \item \emph{el salario de consumo}: con $P_C=P^{1-\alpha}(EP^*)^{\alpha}$, $P_C/P=q^{\alpha}$: una depreciación real encarece la cesta del trabajador frente a lo que vende su empresa.
  \end{itemize}
  En tasas, la OA de corto plazo es la curva de \autor{Phillips} aumentada con expectativas, $\pi_t-\pi^e_t=-a\,(U_t-U^n)+\varepsilon_t$, en la que $\varepsilon_t$ es la perturbación de oferta: la misma que la $u_t$ de la curva neokeynesiana (\ref{sec:nek}), con la letra de cada modelo \vertema{3.A.28}.
\item \emph{A medio plazo}, con las expectativas cumplidas, la oferta es vertical en la producción natural, $Y^n$, la que corresponde a la NAIRU, $U^n$ (la OA de largo plazo, OA$_{\mathrm{LP}}$ en los gráficos). Con la curva de salarios escrita con el índice de consumo, $W/P_C^e=F(U,Z)\cdot\mathrm{PAT}^e$, y la de precios, $W/P=\mathrm{PAT}/(1+\mu)$ \vertema{3.A.28}, la NAIRU de una economía abierta cumple (anexo~\ref{anexo:oa-da})
  \funcionclave[eq:nairu-abierta]{F(U^n,Z)=\frac{1}{(1+\mu)\,q^{\alpha}}}
  \begin{itemize}
  \item \emph{En palabras}: entre lo que cuesta un trabajador a la empresa y lo que este puede comprar hay una \concepto{cuña}: los impuestos sobre el trabajo, el margen y, en una economía abierta, el precio relativo de lo importado\nocite{LayardNickellJackman1991Unemployment}. Si los trabajadores se resisten a que baje su salario real de consumo, una depreciación real o un encarecimiento del petróleo ($q\uparrow$, $\mu\uparrow$) sube la NAIRU y desplaza la OA de largo plazo a la izquierda.
  \end{itemize}
\item \emph{El equilibrio y las políticas a medio plazo}. El equilibrio de corto plazo está donde se cortan la DA y la OA de corto plazo; el de medio plazo, sobre la OA vertical (gráfico~\ref{fig:da-oa-abierta}). Lo que \autor{Mundell}-\autor{Fleming} decía de las políticas cambia cuando los precios se ajustan:
  \begin{itemize}
  \item \emph{con tipo flexible, el dinero es neutral a medio plazo}: una expansión monetaria sube la renta a corto plazo por la depreciación; con $Y>Y^n$, los salarios y los precios suben, y a medio plazo $P$ y $E$ han subido en la misma proporción que $M$: el tipo de cambio real y la renta vuelven a su nivel;
  \item \emph{con tipo fijo, la política fiscal desplaza exportaciones netas}: la expansión sube la renta a corto plazo; a medio plazo, los precios suben, el tipo de cambio real se aprecia y la economía vuelve a $Y^n$ con unas exportaciones netas menores en lo mismo que ha subido el gasto (gráfico~\ref{fig:da-oa-abierta}, panel a). La expulsión que el tipo fijo evitaba a corto plazo llega por los precios;
  \item \emph{sin tipo de cambio propio, el ajuste es lento}: el tipo real solo cambia si los precios nacionales cambian frente a los de fuera. Si $Y<Y^n$, la inflación nacional queda por debajo de la exterior, $q$ sube poco a poco y las exportaciones netas devuelven la economía a $Y^n$ (anexo~\ref{anexo:oa-da}). Es la \concepto{devaluación interna} (\emph{internal devaluation}): tarda y destruye empleo, porque los salarios nominales casi nunca bajan\nocite{SchmittGroheUribe2016Downward} \vertema{3.B.16};
  \item \emph{una devaluación} consigue de una vez lo que el ajuste de precios logra despacio: desplaza la DA y lleva la economía a $Y^n$ sin esperar a los salarios (panel b, de $A$ a $B$). Pero si la economía ya está en $Y^n$, los precios acaban subiendo en la misma proporción y no queda nada, salvo más inflación y menos credibilidad. Fijar una paridad equivocada sale caro: fue la crítica de \textcite{Keynes1925Churchill} a la vuelta del Reino Unido al patrón oro en 1925 con la paridad de antes de la guerra, que dejó la libra sobrevalorada y obligaba a bajar salarios.
  \end{itemize}
\item Volvamos a España: la devaluación interna fue su ajuste tras la crisis. El déficit por cuenta corriente, que llegó al 9,4~\% del PIB en 2007, dio paso a un superávit desde 2012\nocite{BdE2026BE17}, con una caída de la demanda interna, unas exportaciones que crecieron y una moderación salarial \vertema{4.A.21}.
\end{itemize}

\begin{figure}[H]
\centering
\caption{La demanda y la oferta agregadas en una economía abierta con tipo de cambio fijo}\label{fig:da-oa-abierta}
\grafico[0.85\textwidth]{da-oa-abierta}
\fuente{Elaboración propia a partir de \textcite{Blanchard2021Macroeconomics} y \textcite{CarlinSoskice2015Macroeconomics}. Panel a: una expansión fiscal sube la renta a corto plazo (de $A$ a $B$); a medio plazo, la subida de los precios aprecia el tipo real y la economía vuelve a $Y^n$ con menos exportaciones netas ($C$). Panel b: desde una recesión ($A$), la devaluación interna baja los precios, deprecia el tipo real y lleva la economía a $Y^n$ ($C$), despacio; una devaluación ($E\uparrow$) desplaza la DA y la lleva a $Y^n$ de una vez ($B$), con precios más altos.}
\end{figure}

\begin{itemize}
\item \emph{La inflación en una economía abierta}: el índice de consumo mezcla lo nacional y lo importado. En tasas, con $\pi_{H,t}$ la inflación de los bienes nacionales, $\gamma_{E,t}$ la depreciación nominal y $\gamma_{P^*,t}$ la inflación exterior,
  \funcionclave[eq:inflacion-abierta]{\pi_t\simeq(1-\alpha)\,\pi_{H,t}+\alpha\,\big(\gamma_{E,t}+\gamma_{P^*,t}\big)}
  \begin{itemize}
  \item \emph{En palabras}: el segundo sumando es la \concepto{inflación importada}. Una depreciación o un encarecimiento de lo importado suben los precios de consumo aunque los nacionales no cambien: es el \emph{efecto directo}. Si los salarios y los márgenes reaccionan, vienen los \emph{efectos indirectos y de segunda ronda}, que llevan la perturbación a $\pi_{H,t}$ \vertema{3.A.41}.
  \end{itemize}
\end{itemize}

\subparagraph{iii) Implicaciones}

\begin{itemize}
\item \emph{El régimen decide el camino, no el destino}: a medio plazo, la producción vuelve a $Y^n$ y el tipo real es el mismo con los 2 regímenes, pero con tipo fijo el camino puede ser largo y doloroso\nocite{Blanchard2021Macroeconomics}.
\item \emph{Las políticas de demanda no mueven $Y^n$}: a medio plazo, la monetaria solo deja precios y la fiscal cambia la composición de la demanda. Solo lo que desplaza la OA de largo plazo (la cuña, la productividad) cambia la producción de equilibrio: es el terreno de las perturbaciones y de las políticas de oferta (\ref{sec:perturbaciones}).
\item \emph{El exterior entra 2 veces en los precios}: en la oferta, por la cuña, y en el índice de consumo, por la inflación importada. La cuña resta además eficacia a la devaluación: el tipo real más depreciado encarece el consumo y, si los salarios se resisten, rebaja la producción de equilibrio\nocite{CarlinSoskice2015Macroeconomics}.
\item \emph{La economía cerrada es un caso particular}: con $\alpha=0$ desaparece la cuña, y sin movilidad y con tipo flexible la DA es la cerrada.
\end{itemize}

\paragraph{c) Extensiones}

\begin{itemize}
\item \emph{El modelo de demandas rivales}: en el plano $(Y,q)$, la NAIRU abierta es una curva decreciente, la de \emph{equilibrio de las demandas rivales} (EDR) de trabajadores, empresas y proveedores extranjeros sobre la renta; con la curva de equilibrio comercial y la del mercado de bienes, permite seguir a la vez el empleo, la inflación y el saldo exterior tras un encarecimiento de la energía (anexo~\ref{anexo:oa-da})\nocite{LayardNickellJackman1991Unemployment,CarlinSoskice2015Macroeconomics}.
\item \emph{La estanflación con salarios reales rígidos}: \textcite{BrunoSachs1985Economics} explican la de la Europa de los 70 por unos salarios reales que no aceptaron el encarecimiento de las importaciones, que se convirtió en paro.
\end{itemize}

\paragraph{d) Evidencia empírica}

\begin{itemize}
\item \emph{El traspaso} (\emph{pass-through}) del tipo de cambio a los precios es menor y más lento de lo que dice la identidad: muchas importaciones se facturan en dólares y los importadores absorben parte en sus márgenes\nocite{Gopinath2020DCP} \vertema{3.A.37}. En la zona del euro, una depreciación del 1~\% sube los precios de importación un 0,3~\% en un año, pero el IAPC solo un 0,04~\%\nocite{OrtegaOsbat2020ERPT}, y el traspaso es mayor cuando la depreciación nace de una perturbación monetaria\nocite{ComunaleKunovac2017ERPT}.
\item \emph{La relación real de intercambio en 2021-2022}: en 2020, el 80~\% del crudo que importaba la zona del euro se facturaba en dólares; con la energía cara y el euro depreciándose desde mediados de 2021, la pérdida de renta real por el empeoramiento de la relación de intercambio fue de unos 1,3 puntos del PIB en el cuarto trimestre de 2021 frente a un año antes\nocite{GunnellaSchuler2022}.
\end{itemize}

\paragraph{e) Valoración}

\begin{itemize}
\item \emph{Lo que aporta}: con 2 curvas determina a la vez la renta y los precios, distingue el corto del medio plazo y separa lo que hace la demanda de lo que hace la oferta, que es lo que pide el título.
\item \emph{Sus límites}: las expectativas, la curva de \autor{Phillips} y la política monetaria entran sin microfundamentos; trabaja con niveles y no con la dinámica de la inflación, y no tiene un criterio de bienestar para comparar políticas. El paso siguiente es reconstruirlo con agentes que optimizan.
\end{itemize}

% --------------------------------------------------------------------------
\subsection{El modelo neokeynesiano}\label{sec:nek}

\paragraph{a) Idea}

\begin{itemize}
\item La pregunta: \emph{¿cómo se determinan la producción y la inflación cuando los precios son rígidos, pero los agentes optimizan y tienen expectativas racionales?}
\item La respuesta de la nueva economía keynesiana (NEK)\nocite{ClaridaGaliGertler1999Science,Woodford2003Interest,Gali2015Monetary} \vertema{3.A.7} es un modelo de 3 ecuaciones: la IS dinámica hace de demanda agregada, la curva de \autor{Phillips} neokeynesiana hace de oferta agregada y una regla de tipos de interés sustituye a la LM. Es la DA-OA del apartado~\ref{sec:da-oa} con microfundamentos, en inflación y no en niveles de precios, y es el que usan hoy los bancos centrales.
\end{itemize}

\paragraph{b) Modelo}

\subparagraph{i) Supuestos}

\begin{itemize}
\item \emph{Economía cerrada} (se abre en el bloque~\ref{sec:abierta}), sin capital y sin dinero más que como unidad de cuenta (\emph{cashless}).
\item \emph{Hogares} representativos de vida infinita, con aversión relativa al riesgo $\theta$ (la inversa de la elasticidad de sustitución intertemporal) y desutilidad del trabajo con la inversa de la elasticidad de \autor{Frisch}, $\varphi$.
\item \emph{Empresas} en \emph{competencia monopolística} à la \textcite{DixitStiglitz1977Monopolistic},\footnote{\autor{Joseph Stiglitz} recibió el Premio Nobel de Economía en 2001, con \autor{George Akerlof} y \autor{Michael Spence}, «por sus análisis de los mercados con información asimétrica»\nocite{Nobel2001AkerlofSpenceStiglitz}.} con elasticidad de sustitución $\sigma$ \vertema{3.A.18}, que producen con trabajo y una tecnología $A_t$.
\item \emph{Precios rígidos} à la \textcite{Calvo1983Staggered}: cada periodo, una empresa no revisa su precio con probabilidad $\eta$.
\item \emph{4 perturbaciones}: tecnológica ($A_t$), de demanda (preferencias o gasto público), de costes ($u_t$) y de política monetaria ($\varepsilon_t$). Las ecuaciones se log-linealizan alrededor de un estado estacionario con inflación cero; $x_t\equiv y_t-y^n_t$ es la \concepto{brecha de producción} (\emph{output gap}), la distancia a la producción que habría con precios flexibles.
\end{itemize}

\subparagraph{ii) Desarrollo}

\begin{itemize}
\item \emph{La demanda: la curva IS dinámica}. De la ecuación de \autor{Euler} del hogar, log-linealizada \vertema{3.A.29}:
  \funcionclave[eq:is-dinamica]{x_t=\mathbb{E}_t[x_{t+1}]-\frac{1}{\theta}\Big(i_t-\mathbb{E}_t[\pi_{t+1}]-r^n_t\Big)}
  \begin{itemize}
  \item \emph{En palabras}: si el tipo real supera al \concepto{tipo de interés natural}, $r^n_t$ (el que habría con precios flexibles), los hogares aplazan consumo y la brecha cae. Las perturbaciones de demanda (preferencias, gasto público) entran subiendo $r^n_t$; una mejora tecnológica transitoria lo baja, porque se espera que la producción natural vuelva a caer.
  \item Es \emph{prospectiva}: la brecha de hoy depende de toda la senda esperada de tipos reales. Por eso importan los anuncios del banco central.
  \end{itemize}
\item \emph{La oferta: la curva de \autor{Phillips} neokeynesiana} (NKPC). Las empresas que revisan su precio lo fijan con un margen sobre los costes marginales que esperan mientras siga vigente:
  \funcionclave[eq:nkpc]{\pi_t=\beta\,\mathbb{E}_t[\pi_{t+1}]+\kappa\,x_t+u_t,\qquad \kappa\equiv\frac{(1-\eta)(1-\beta\eta)}{\eta}\,(\theta+\varphi)}
  \begin{itemize}
  \item \emph{En palabras}: la inflación depende de la esperada y de la brecha, que mide el coste marginal real; cuanto más rígidos los precios, menor $\kappa$ y más plana la curva. $u_t$ es la \emph{perturbación de costes} (márgenes deseados, impuestos, salarios) que no viene de la brecha: es la $\varepsilon_t$ de la OA del apartado~\ref{sec:da-oa}.
  \end{itemize}
\item \emph{La política monetaria: la regla de \autor{Taylor}} \vertema{3.A.36}, que sustituye a la LM:
  \begin{equation}
  i_t=\rho+\phi_\pi\,\pi_t+\phi_x\,x_t+\varepsilon_t,\qquad \rho\equiv-\ln\beta, \label{eq:taylor}
  \end{equation}
  con $\varepsilon_t$ lo que hace el banco central al margen de la regla. \textcite{Taylor1993} la propuso como una descripción de la Reserva Federal, con $\phi_\pi=1{,}5$ y $\phi_x=0{,}5$.
\item \emph{La DA y la OA neokeynesianas en el plano $(x_t,\pi_t)$}. Con expectativas dadas (perturbaciones transitorias), quedan 2 rectas (gráfico~\ref{fig:nek-perturbaciones}; anexo~\ref{anexo:nek}):
  \begin{itemize}
  \item la \emph{DA} (la IS dinámica con la regla dentro), decreciente: más inflación sube el tipo de interés y reduce la brecha. La desplazan el tipo natural ($r^n_t$) y la política monetaria ($\varepsilon_t$):
    \begin{equation}
    \pi_t=-\frac{\theta+\phi_x}{\phi_\pi}\,x_t+\frac{r^n_t-\rho-\varepsilon_t}{\phi_\pi} ; \label{eq:da-nek}
    \end{equation}
  \item la \emph{OA}, la NKPC, creciente con pendiente $\kappa$; la desplaza $u_t$.
  \end{itemize}
\end{itemize}

\begin{figure}[H]
\centering
\caption{Las perturbaciones y la política en el modelo neokeynesiano, en el plano $(x_t,\pi_t)$}\label{fig:nek-perturbaciones}
\grafico[0.8\textwidth]{nek-perturbaciones}
\fuente{Elaboración propia a partir de \textcite{Gali2015Monetary} y \textcite{ClaridaGaliGertler1999Science}, con perturbaciones transitorias. Panel a: una mejora tecnológica transitoria baja el tipo natural y la DA; si el tipo de interés sigue al natural, la producción sube hasta la natural y la economía vuelve a $E_0$. Panel b: una perturbación de costes sube la NKPC; el banco central solo puede elegir entre $E_\pi$ (inflación en el objetivo, más recesión), $E_x$ (brecha cerrada, más inflación) o puntos intermedios como $E_1$. Panel c: una política de oferta baja la NKPC y mejora las 2 ($E_2$). Panel d: una perturbación de demanda mueve la brecha y la inflación en el mismo sentido; si el tipo de interés sigue al natural, vuelve a $E_0$ (divina coincidencia).}
\end{figure}

\begin{anotaciones}
\textbf{Por qué el modelo neokeynesiano se puede dibujar en el plano $(x_t,\pi_t)$, y con qué supuestos}
\begin{itemize}
\item \emph{En sentido estricto, no se puede}: el modelo canónico es un sistema de ecuaciones en diferencias estocásticas con expectativas, no un par de curvas estáticas. La NKPC y la IS contienen $\mathbb{E}_t[\pi_{t+1}]$ y $\mathbb{E}_t[x_{t+1}]$, que son endógenas, y no hay una «curva de oferta» ni una «de demanda» que se desplacen.
\item \emph{Los supuestos que hay detrás del gráfico}:
  \begin{itemize}
  \item las ecuaciones están \emph{linealizadas} alrededor de un \emph{estado estacionario} sin inflación ni brecha, y solo valen para desviaciones pequeñas;
  \item las perturbaciones siguen procesos \emph{AR(1)}, por ejemplo $u_t=\rho_u\,u_{t-1}+\epsilon^u_t$, con $0\leq\rho_u<1$ ($\rho_u$ y $\epsilon^u_t$, locales);
  \item se busca la solución de \emph{variables de estado mínimas}, de la forma $x_t=\psi_x\,u_t$ y $\pi_t=\psi_\pi\,u_t$ (con $\psi_x$ y $\psi_\pi$ por determinar);
  \item de ahí, \emph{expectativas proporcionales}: $\mathbb{E}_t[x_{t+1}]=\rho_u\,x_t$ y $\mathbb{E}_t[\pi_{t+1}]=\rho_u\,\pi_t$.
  \end{itemize}
\item \emph{De dónde sale la representación}: con esas expectativas, la NKPC queda $\pi_t=\frac{\kappa}{1-\beta\rho_u}\,x_t+\frac{u_t}{1-\beta\rho_u}$, creciente (la «oferta»), y la IS con la regla de \autor{Taylor} (sin $\varepsilon_t$ y con $\hat r^n_t\equiv r^n_t-\rho$),
  \begin{equation*}
  \pi_t=-\frac{(1-\rho_u)\,\theta+\phi_x}{\phi_\pi-\rho_u}\,x_t+\frac{\hat r^n_t}{\phi_\pi-\rho_u},
  \end{equation*}
  decreciente si $\phi_\pi>\rho_u$ (la «demanda»). El gráfico~\ref{fig:nek-perturbaciones} y la ecuación~\eqref{eq:da-nek} son el caso de perturbaciones transitorias, $\rho_u=0$, en el que las expectativas son nulas.
\item \emph{Cómo se lee}: las 2 rectas relacionan las \emph{respuestas de equilibrio} de $x_t$ y $\pi_t$ a una perturbación con una persistencia dada; no son sendas por las que se mueva la economía. Si cambia la persistencia o la regla, cambian las 2 rectas. En la economía abierta (\ref{sec:gali-monacelli}) vale lo mismo con $\pi_{H,t}$, $\theta_\alpha$ y $\kappa_\alpha$.
\end{itemize}
\end{anotaciones}

\begin{itemize}
\item \emph{La forma reducida: los multiplicadores de las 2 políticas}. Resolviendo las 2 rectas (anexo~\ref{anexo:nek}),
  \funcionclave[eq:forma-reducida]{x_t=\frac{\hat r^n_t-\phi_{\pi}\,u_t-\varepsilon_t}{\theta+\phi_x+\phi_{\pi}\,\kappa}}
  \begin{itemize}
  \item \emph{En palabras}: la brecha se mueve por 3 fuentes, el tipo natural, los costes y la política monetaria. El multiplicador de la política monetaria, $-1/(\theta+\phi_x+\phi_\pi\kappa)$, y el de la fiscal (que actúa subiendo $r^n_t$) comparten denominador y son menores cuanto más agresiva es la regla con la inflación: la política monetaria sistemática contrarresta cualquier impulso a la demanda.
  \end{itemize}
\end{itemize}

\subparagraph{iii) Implicaciones}

\begin{itemize}
\item \emph{El dinero no es neutral a corto plazo, pero sí a largo}: con precios rígidos, subir $i_t$ sube el tipo real, baja la brecha y, con ella, la inflación; a largo plazo, con $\beta$ casi 1, la NKPC es casi vertical, como la OA de largo plazo del apartado~\ref{sec:da-oa}.
\item \emph{Demanda frente a costes}: una perturbación de demanda desplaza la DA y mueve la brecha y la inflación en el mismo sentido; una de costes desplaza la NKPC y las mueve en sentido contrario. Es la distinción que ordena el bloque~\ref{sec:perturbaciones}.
\item \emph{La \concepto{divina coincidencia}}: si no hay perturbaciones de costes ($u_t=0$), estabilizar la inflación estabiliza también la brecha ($\pi_t=0\Rightarrow x_t=0$), y basta con que el tipo de interés siga al natural\nocite{BlanchardGali2007Real}.
\item \emph{Un criterio de bienestar}: la pérdida de bienestar depende de la volatilidad de la inflación y de la brecha, de modo que se pueden comparar reglas y regímenes \vertema{3.A.36}.
\end{itemize}

\paragraph{c) Extensiones}

\begin{itemize}
\item \emph{Inercia de la inflación}: con indexación o con empresas que siguen una regla sencilla sale una curva de \autor{Phillips} híbrida, con inflación pasada, que da la persistencia que se ve en los datos\nocite{GaliGertler1999Inflation} \vertema{3.A.41}.
\item \emph{Hogares heterogéneos}: si parte de los hogares consume lo que gana, la política fiscal y la monetaria actúan sobre todo a través de la renta y no del tipo de interés\nocite{GaliLopezSalidoValles2007,KaplanMollViolante2018Monetary} (\ref{sec:demanda-origen}).
\item \emph{La economía abierta}: \textcite{GaliMonacelli2005Monetary} mantienen la estructura con otros coeficientes y 2 inflaciones (\ref{sec:gali-monacelli}).
\end{itemize}

\paragraph{d) Evidencia empírica}

\begin{itemize}
\item \emph{Tecnología y horas}: tras una mejora tecnológica las horas trabajadas \emph{caen}, como predice el modelo con precios rígidos y al contrario que el del ciclo real\nocite{Gali1999Technology,BasuFernaldKimball2006Are} \vertema{3.A.42}.
\item \emph{La curva de \autor{Phillips} es plana}: la inflación respondió poco a la brecha durante décadas, lo que se lee como un $\kappa$ bajo o como expectativas bien ancladas, y se empinó con la energía y los cuellos de botella de 2021-2022\nocite{BernankeBlanchard2025What} \vertema{3.A.41}.
\end{itemize}

\paragraph{e) Valoración}

\begin{itemize}
\item \emph{Lo que aporta}: es estructural y microfundamentado, de modo que la crítica de \autor{Lucas}\nocite{Lucas1976Econometric}\footnote{\autor{Robert Lucas} recibió el Premio Nobel de Economía en 1995 «por haber desarrollado y aplicado la hipótesis de las expectativas racionales, transformando así el análisis macroeconómico y profundizando en la comprensión de la política económica»\nocite{Nobel1995Lucas}.} no le alcanza \vertema{3.A.6}; separa con claridad las perturbaciones de demanda de las de costes y da un criterio de bienestar para la política.
\item \emph{Sus límites}: necesita inercia añadida para reproducir la inflación; la $u_t$ es poco más que un residuo; la versión básica no tiene sector financiero, y la orientación futura (\emph{forward guidance}) tiene en él efectos poco creíbles.
\item \emph{Recapitulación de I}: \autor{Mundell}-\autor{Fleming} decide qué política funciona según el régimen; la DA-OA abierta añade los precios, la cuña exterior y la inflación importada; el modelo neokeynesiano da la versión moderna de las 2 curvas. Con este marco, pasamos a lo que ocurre cuando una perturbación saca a la economía del equilibrio.
\end{itemize}


% ==========================================================================
\section{Las perturbaciones de oferta y de demanda}\label{sec:perturbaciones}
\minutos{10}

\begin{itemize}
\item Con el marco construido, seguimos cada perturbación desde su origen: cómo se propaga, qué efectos de segunda ronda tiene y cómo debe responder la política económica. Empezamos por las de oferta, que son las que pusieron a prueba el modelo en los 70 y en 2022, y seguimos con las de demanda, que a menudo nacen de ellas.
\item \emph{Conclusión anticipada}:
  \begin{itemize}
  \item las perturbaciones de oferta mueven la renta y los precios en sentido contrario; las tecnológicas no plantean dilema a la política monetaria, pero las de costes sí, y la política de demanda solo puede repartir el golpe; las políticas de oferta son las únicas que mejoran las 2 a la vez;
  \item las de demanda mueven la renta y los precios en el mismo sentido y la política monetaria puede estabilizarlas sin coste (divina coincidencia); cuando no puede, porque los tipos están en su límite inferior o porque no hay moneda propia, queda la fiscal;
  \item las 2 se mezclan: una perturbación de oferta puede acabar siendo de demanda, y una de demanda puede chocar con una oferta que no responde.
  \end{itemize}
\end{itemize}

% --------------------------------------------------------------------------
\subsection{Las perturbaciones de oferta: origen y propagación}\label{sec:oferta-origen}

\begin{itemize}
\item \emph{Definición}: una \concepto{perturbación de oferta} cambia, a cada nivel de precios, lo que las empresas quieren producir: desplaza la OA. Por eso mueve la producción y los precios en sentido contrario.
\item \emph{Su origen}, en 3 tipos:
  \begin{itemize}
  \item \emph{tecnológicas}: cambios de la productividad, que mueven la producción natural; son las del ciclo real \vertema{3.A.42}. En el modelo neokeynesiano, una mejora transitoria baja el tipo natural: la producción sube, pero menos que la natural, y la brecha y la inflación caen (gráfico~\ref{fig:nek-perturbaciones}, panel a);
  \item \emph{de costes y márgenes}: salarios, márgenes o impuestos indirectos que suben sin que cambie la productividad; son la $u_t$ de la NKPC, que sube: menos brecha y más inflación (panel b) \vertema{3.A.41};
  \item \emph{externas}: el precio de la energía y de las materias primas, los cortes de las cadenas de suministro o el clima. En una economía abierta son las más frecuentes, y suelen ser a la vez de oferta (encarecen producir) y de demanda (transfieren renta a los exportadores de energía).
  \end{itemize}
\item Volvamos a España: importa en términos netos el 69~\% de la energía que consume, frente al 57~\% de la UE\nocite{Eurostat2026nrgindid} \vertema{4.A.6}. Un encarecimiento de la energía es para ella una perturbación de oferta externa, y en la DA-OA abierta tiene 3 efectos (gráfico~\ref{fig:oferta-abierta}, panel a):
  \begin{itemize}
  \item \emph{la OA de corto plazo sube}: producir cuesta más, y la inflación de consumo sube además por el efecto directo de la energía \eqref{eq:inflacion-abierta};
  \item \emph{la OA de largo plazo se desplaza a la izquierda} si los salarios reales se resisten a bajar: la cuña exterior sube la NAIRU \eqref{eq:nairu-abierta};
  \item \emph{la DA cae}: empeora la relación real de intercambio y parte de la renta se va a pagar la energía de fuera.
  \end{itemize}
  El resultado es la \concepto{estanflación}: menos producción y más inflación a la vez.
\end{itemize}

\begin{figure}[H]
\centering
\caption{Una perturbación de oferta adversa y una política de oferta en una economía abierta}\label{fig:oferta-abierta}
\grafico[0.85\textwidth]{oferta-abierta}
\fuente{Elaboración propia a partir de \textcite{Blanchard2021Macroeconomics} y \textcite{BrunoSachs1985Economics}. Panel a: el encarecimiento de la energía importada sube la OA de corto plazo, desplaza a la izquierda la de largo plazo (si los salarios reales se resisten) y reduce algo la DA: menos producción y más precios ($A\to B$). Panel b: una política de oferta desplaza las 2 OA a la derecha: más producción a corto plazo ($B$) y a medio plazo ($C$), con precios más bajos y un tipo de cambio real más depreciado.}
\end{figure}

\begin{itemize}
\item \emph{Su propagación: del efecto directo a la segunda ronda}.
  \begin{itemize}
  \item El \emph{efecto directo} es la subida del precio de lo que se consume directamente (los carburantes, la electricidad); el \emph{indirecto}, la de los bienes que usan energía para producirse; la \emph{segunda ronda}, la de los salarios y los márgenes que intentan recuperar lo perdido. Las 2 primeras son un cambio de precios relativos y pasan; la tercera convierte la perturbación en inflación persistente \vertema{3.A.41}.
  \item Un cambio de precios relativos sube el nivel general porque, con costes de cambiar los precios, suben antes los pocos precios que quieren subir mucho que los muchos que quieren bajar un poco\nocite{BallMankiw1995Relative}.
  \end{itemize}
\item \emph{La heterogeneidad sectorial: las perturbaciones no se diluyen}.
  \begin{itemize}
  \item Si las empresas grandes pesan mucho, lo que les pasa no se compensa en el agregado: es la hipótesis \emph{granular}\nocite{Gabaix2011Granular}.
  \item Por la red de compras y ventas entre sectores, una perturbación de oferta en uno viaja \emph{aguas abajo}, hacia sus clientes, que lo usan como insumo, con un efecto total mayor que el directo\nocite{AcemogluAkcigitKerr2016} (anexo~\ref{anexo:redes}).\footnote{\autor{Daron Acemoglu} recibió el Premio Nobel de Economía en 2024, con \autor{Simon Johnson} y \autor{James Robinson}, «por sus estudios sobre cómo se forman las instituciones y cómo afectan a la prosperidad»\nocite{NobelPrize2024}.} Así, los semiconductores en 2021 o la energía en 2022 movieron la economía mucho más de lo que pesan en el PIB \vertema{3.A.31}.
  \end{itemize}
\item \emph{La identificación}: las perturbaciones no se observan y llegan mezcladas. Los vectores autorregresivos estructurales (\emph{structural VAR}) de \textcite{Sims1980Macroeconomics}\footnote{\autor{Christopher Sims} recibió el Premio Nobel de Economía en 2011, con \autor{Thomas Sargent}, «por su investigación empírica sobre la causa y el efecto en la macroeconomía»\nocite{Nobel2011SargentSims}.} las separan con restricciones que vienen de la teoría: \textcite{BlanchardQuah1989} suponen que solo las de oferta tienen efectos permanentes sobre la producción, y \textcite{Kilian2009NotAll} muestra que no es lo mismo un corte de la oferta de crudo que un aumento de la demanda mundial \vertema{3.A.37}.
\end{itemize}

\begin{anotaciones}
\textbf{Cómo se identifican las perturbaciones}
\begin{itemize}
\item \emph{\autor{Blanchard} y \autor{Quah}}: con la producción y el paro de Estados Unidos, la restricción es que la perturbación de demanda, $\varepsilon^{\mathrm D}_t$, no tiene efecto a largo plazo sobre la producción ($\varepsilon^{\mathrm D}_t$ y $\varepsilon^{\mathrm S}_t$, con superíndice en redonda, son locales de esta caja):
  \begin{equation}
  \lim_{k\to\infty}\frac{\partial y_{t+k}}{\partial\varepsilon^{\mathrm D}_t}=0 . \label{eq:blanchard-quah}
  \end{equation}
  El efecto de la demanda sobre la producción y el paro tiene forma de joroba y desaparece en 2 a 5 años; el de la oferta, $\varepsilon^{\mathrm S}_t$, crece con el tiempo y se estabiliza a los 5. Una perturbación de oferta favorable puede subir algo el paro al principio, y después lo baja y vuelve despacio a su nivel (gráfico~\ref{fig:blanchard-quah}). Cuánto pesa cada una depende de los supuestos: la demanda explica entre el 40 y más del 95~\% de las fluctuaciones a un año\nocite{BlanchardQuah1989}.
\begin{figuracaja}
\centering
\caption{Las respuestas a una perturbación de demanda y a una de oferta en el modelo de \autor{Blanchard} y \autor{Quah}}\label{fig:blanchard-quah}
\grafico[0.85\textwidth]{blanchard-quah}
\fuente{Elaboración propia a partir de \textcite{BlanchardQuah1989}. Esquema cualitativo, no las estimaciones. La demanda mueve la producción y el paro en sentidos opuestos, en joroba, y su efecto desaparece; la oferta deja la producción en un nivel permanentemente más alto y no afecta al paro a largo plazo.}
\end{figuracaja}
\item \emph{\autor{Galí}}: la tecnología es la única perturbación con efectos permanentes sobre la productividad del trabajo; tras una mejora, las horas \emph{caen}, y la tecnología explica poco del ciclo\nocite{Gali1999Technology}.
\item \emph{\autor{Kilian}}: separa las perturbaciones de la oferta de crudo, de la demanda mundial de materias primas y de la demanda específica de petróleo; cada una tiene efectos distintos, y la subida de los 2000 se debió sobre todo a la demanda mundial\nocite{Kilian2009NotAll}.
\item \emph{\autor{Hamilton}}: todas las recesiones de Estados Unidos desde la Segunda Guerra Mundial salvo una vinieron precedidas de una fuerte subida del petróleo\nocite{Hamilton1983Oil}.
\end{itemize}
\end{anotaciones}

% --------------------------------------------------------------------------
\subsection{La respuesta a las perturbaciones de oferta}\label{sec:oferta-respuesta}

\begin{itemize}
\item \emph{Las tecnológicas no plantean dilema}: si el tipo de interés sigue al natural, la producción sube hasta la natural y la inflación no se mueve (gráfico~\ref{fig:nek-perturbaciones}, panel a). Para la política monetaria se tratan como las de demanda: lo que importa no es si mueven la oferta, sino si mueven la NKPC.
\item \emph{Las de costes sí: el dilema}. Contraer la demanda frena la inflación pero agrava la caída de la producción; acomodar la sostiene, pero arriesga los efectos de segunda ronda \vertema{3.A.41}.
  \begin{itemize}
  \item En el modelo neokeynesiano no hay divina coincidencia: el banco central solo puede elegir un punto de la NKPC nueva, entre la inflación en el objetivo con más recesión ($E_\pi$) y la brecha cerrada con más inflación ($E_x$); con su regla, reparte el golpe ($E_1$; panel b) \vertema{3.A.36}. Cuanto más pesan las perturbaciones de costes y más rígidos son los precios, menos agresiva debe ser la respuesta a la inflación\nocite{Gali2015Monetary}.
  \item La coincidencia se rompe también sin perturbaciones de costes si los salarios reales son rígidos: estabilizar la inflación ya no estabiliza la brecha que importa para el bienestar\nocite{BlanchardGali2007Real}.
  \end{itemize}
\item \emph{Cuándo mirar a través} (\emph{look through}): ante una perturbación pasajera y con las expectativas bien ancladas, el banco central puede dejar pasar el efecto directo; si es persistente o contagia a salarios y márgenes, tiene que responder, y cuanto antes, mejor\nocite{BCE2025Evaluacion}.
\item \emph{La política fiscal}: con 2 objetivos (renta e inflación) y una perturbación de costes, la monetaria no basta.
  \begin{itemize}
  \item Las ayudas deben ser \emph{temporales y focalizadas} en los hogares más vulnerables: las generalizadas a la energía sostienen la demanda que la política monetaria quiere frenar. Las medidas europeas bajaron la inflación de 2022 en 1 o 2 puntos, pero la mayoría no se dirigió a quien más lo necesitaba\nocite{DaoEtAl2023Unconventional}.
  \item Las medidas sobre el precio, como el tope al gas de la generación eléctrica en España y Portugal, bajan la inflación medida, pero no la cantidad de energía que se demanda.
  \end{itemize}
\end{itemize}

\begin{anotaciones}
\textbf{España ante la energía cara de 2022}
\begin{itemize}
\item La «excepción ibérica» topó desde el 14 de junio de 2022 el precio del gas de las centrales que fijan el precio de la electricidad: 40 euros por megavatio hora durante 6 meses y después una senda creciente, hasta 65 euros en diciembre de 2023, cuando terminó\nocite{BOE2022RDL10,BOE2022OrdenTED517,BOE2023RDL3,CE2022AyudaIberica}.
\item España bajó el IVA de la electricidad del 21~\% al 10~\% en 2021 y al 5~\% en la segunda mitad de 2022\nocite{BOE2021RDL12,BOE2022RDL11}, y subvencionó con 20 céntimos cada litro de carburante de abril a diciembre de 2022\nocite{BOE2022RDL6,BOE2022RDL11}: medidas generales, no focalizadas.
\item En España la energía llegó a subir un 60,3~\% en marzo de 2022, pero en octubre solo un 7,9~\%, frente al 41,5~\% de la zona del euro\nocite{Eurostat2026IAPC}: había subido antes y se sumaron el tope al gas y la rebaja del IVA.
\end{itemize}
\end{anotaciones}

\begin{itemize}
\item \emph{Las políticas de oferta: atacar el origen}. Una \concepto{política de oferta} aumenta la capacidad productiva de la economía o reduce sus costes. Actúa en 3 frentes:
  \begin{itemize}
  \item \emph{los mercados}: la competencia en los mercados de bienes baja los márgenes; las reformas del mercado de trabajo, la cuña fiscal y las prestaciones bajan la NAIRU \vertema{3.A.28};
  \item \emph{la productividad}: el capital humano, la innovación, las infraestructuras y la calidad de las instituciones suben la producción potencial \vertema{3.A.45};
  \item \emph{la energía}: diversificar las fuentes y aumentar las renovables reduce la exposición a las perturbaciones de fuera \vertema{4.A.6}.
  \end{itemize}
\item \emph{Sus efectos}: en el modelo neokeynesiano, bajar los márgenes o los costes reduce $u_t$ y baja la NKPC: la brecha y la inflación mejoran a la vez, sin dilema, porque la política actúa sobre el origen de la perturbación (gráfico~\ref{fig:nek-perturbaciones}, panel c). En la DA-OA abierta desplaza a la derecha las 2 OA: más producción y precios más bajos, que con el tipo de cambio dado deprecian el tipo real y ganan competitividad (gráfico~\ref{fig:oferta-abierta}, panel b).
\item \emph{Las políticas de rentas}: los acuerdos sobre salarios y precios con que se quiso frenar la estanflación de los 70 buscaban lo mismo por otra vía: bajar la cuña sin pasar por el paro \vertema{3.A.5}.
  \begin{itemize}
  \item Un pacto de rentas entre el Gobierno, los sindicatos y las empresas modera los salarios que se piden y reparte la caída del salario real que impone la perturbación: menos paro y menos inflación. Fue la vía de los países escandinavos y de Austria tras el petróleo\nocite{CarlinSoskice2015Macroeconomics}.
  \item Que funcione depende de cómo se negocia: con negociación muy centralizada, los sindicatos tienen en cuenta el efecto de sus salarios sobre los precios y el paro; muy descentralizada, la competencia contiene los salarios; por sectores, ni lo uno ni lo otro, y el salario real y el paro son los más altos\nocite{CalmforsDriffill1988}.\footnote{La relación entre la centralización de la negociación colectiva y el salario real o el paro tiene forma de joroba: bajos en los 2 extremos y máximos en el nivel intermedio.\graficonota[0.5]{calmfors-driffill}}
  \end{itemize}
\item \emph{Lo que dice la evidencia}:
  \begin{itemize}
  \item \emph{a largo plazo, ganan}: desregular los mercados de bienes reduce las rentas y sube el empleo y el salario real; desregular el de trabajo redistribuye las rentas y tiene costes a corto plazo para los trabajadores\nocite{BlanchardGiavazzi2003};
  \item \emph{a corto plazo, depende del ciclo}: las de la protección del empleo y las prestaciones por desempleo ayudan en las expansiones y pueden ser contractivas cuando hay holgura, y las rebajas de la cuña fiscal y las políticas activas de empleo funcionan mejor en las recesiones\nocite{DuvalFurceri2018};
  \item \emph{con los tipos en el límite inferior}, bajar los márgenes reduce la inflación esperada, sube el tipo real y puede contraer la demanda a corto plazo\nocite{EggertssonFerreroRaffo2014};
  \item \emph{en España}, las reformas laborales de 2012\nocite{BOE2012RDL3,BOE2012Ley3} y de 2021\nocite{BOE2021RDL32} son ejemplos de políticas de oferta en el mercado de trabajo; la OCDE atribuyó a la de 2012 una contribución a la moderación salarial\nocite{OCDE2014ReformaLaboral} \vertema{4.A.16}.
  \end{itemize}
\item \emph{En la agenda actual}: el informe \autor{Draghi} sitúa la brecha de productividad con Estados Unidos en el centro del problema europeo\nocite{CE2024Draghi}; el Mecanismo de Recuperación y Resiliencia condiciona sus fondos a reformas e inversiones\nocite{UE2021Reglamento241}, y la OCDE ha sustituido su informe \emph{Going for Growth} por uno nuevo sobre los fundamentos del crecimiento\nocite{OCDE2026Fundamentos}. No son la «economía de la oferta» (\emph{supply-side economics}) de los 80, que esperaba que bajar los impuestos se pagara solo, sino reformas lentas, con perdedores, que necesitan una demanda que no se hunda mientras dan fruto.
\item \emph{2 episodios}:
  \begin{itemize}
  \item \emph{Los 70}: tras el embargo de los países árabes de la OPEP de 1973 y la revolución iraní de 1979, el precio del petróleo se multiplicó casi por 4 entre 1973 y 1974 y casi se triplicó entre 1978 y 1980\nocite{Hamilton2011Historical,BancoMundial2026PinkSheet}. Las economías europeas, con salarios indexados y una política monetaria que acomodó, tuvieron a la vez un paro creciente e inflación de 2 cifras\nocite{BrunoSachs1985Economics}; estimular la demanda cuando el equilibrio con inflación estable ya exigía más paro solo sirvió para que la inflación siguiera subiendo\nocite{CarlinSoskice2015Macroeconomics}. En España, más dependiente del petróleo y con salarios indiciados, la inflación llegó al 28,4~\% en agosto de 1977\nocite{INE2026IPCSerie} \vertema{4.A.27}.
  \item \emph{Los 2000}: el petróleo subió con fuerza, pero sin estanflación. \textcite{BlanchardGali2010Oil} lo explican por 4 razones: hubo menos perturbaciones a la vez, el petróleo pesaba menos en la producción, los salarios reales eran más flexibles y la política monetaria era más creíble.\footnote{Es la versión petrolera del debate sobre la Gran Moderación: de la caída de la volatilidad del producto en Estados Unidos, entre el 20 y el 30~\% se debería a una mejor política, otro 20-30~\% a la buena suerte identificable (perturbaciones de productividad y de materias primas más pequeñas) y el 40-60~\% restante a perturbaciones más pequeñas sin identificar\nocite{StockWatson2003Has}.} Además, aquella subida la empujó sobre todo la demanda mundial, no un corte de la oferta\nocite{Kilian2009NotAll}.
  \end{itemize}
\item La lección de los 2 episodios: con una perturbación de oferta, la política de demanda decide cuánto dura el daño, y la de oferta, cuánto queda. Pero no todas las perturbaciones de oferta se quedan en la oferta.
\end{itemize}

\begin{anotaciones}
\textbf{Las cifras de los 70}
\begin{itemize}
\item \emph{1973}: el 17 de octubre, en plena guerra del Yom Kipur, los países árabes de la OPEP anunciaron un embargo; en noviembre producían 4,4 millones de barriles diarios menos que en septiembre, el 7,5~\% de la producción mundial, y el 1 de enero de 1974 los países del golfo Pérsico duplicaron el precio\nocite{Hamilton2011Historical}. El precio medio del crudo pasó de 2,8 a 11,0 dólares por barril entre 1973 y 1974\nocite{BancoMundial2026PinkSheet}.
\item \emph{1979-1980}: la revolución iraní restó 4,8 millones de barriles diarios, el 7~\% de la producción mundial, entre octubre de 1978 y enero de 1979, y la guerra entre Irán e Irak, desde septiembre de 1980, otro 6~\%\nocite{Hamilton2011Historical}; el precio pasó de 12,9 dólares por barril en 1978 a 36,9 en 1980\nocite{BancoMundial2026PinkSheet}.
\item \emph{Los efectos}: en la OCDE la inflación llegó al 15,5~\% en 1974 y al 13,6~\% en 1980, y el crecimiento cayó al 0,4~\% en 1975 y al 0,2~\% en 1982\nocite{BancoMundial2026WDIhist}. España apenas creció en 1975 (0,5~\%) y en 1979 (0~\%)\nocite{CE2026AMECO}.
\end{itemize}
\end{anotaciones}

% --------------------------------------------------------------------------
\subsection{Las perturbaciones de demanda y su interacción con las de oferta}\label{sec:demanda-origen}

\begin{itemize}
\item \emph{Definición}: una \concepto{perturbación de demanda} cambia, a cada nivel de precios, lo que se quiere gastar: desplaza la DA y mueve la producción y los precios en el mismo sentido.
\item \emph{Su origen}: en el modelo neokeynesiano, casi todas entran subiendo o bajando el tipo natural (gráfico~\ref{fig:nek-perturbaciones}, panel d):
  \begin{itemize}
  \item \emph{preferencias y confianza}: un aumento del ahorro por precaución, como tras la crisis de 2008, baja $r^n_t$;
  \item \emph{la política fiscal}: más gasto público sube $r^n_t$;
  \item \emph{el crédito y la riqueza}: una caída del precio de los activos o un endurecimiento del crédito reducen el gasto;
  \item \emph{la política monetaria}: una subida de tipos al margen de la regla ($\varepsilon_t>0$) desplaza la DA a la izquierda;
  \item \emph{la demanda exterior}, en una economía abierta (\ref{sec:externas}).
  \end{itemize}
\item \emph{Su propagación}:
  \begin{itemize}
  \item \emph{el multiplicador}: el gasto de unos es la renta de otros, y en la DA-OA el efecto se reparte entre producción y precios según la pendiente de la OA;
  \item \emph{el acelerador financiero}: la prima de financiación externa sube cuando cae el patrimonio de los prestatarios, de modo que una perturbación que reduce ese patrimonio encarece el crédito y deprime aún más la inversión\nocite{BernankeGertlerGilchrist1999};\footnote{\autor{Ben Bernanke} recibió el Premio Nobel de Economía en 2022, con \autor{Douglas Diamond} y \autor{Philip Dybvig}, «por sus investigaciones sobre los bancos y las crisis financieras»\nocite{Nobel2022Bancos}.}
  \item \emph{los hogares que viven al día}: si una parte de los hogares consume lo que gana, el consumo sube tras un aumento del gasto público, en vez de caer como predice la equivalencia ricardiana\nocite{GaliLopezSalidoValles2007} \vertema{3.A.33};
  \item \emph{la red de producción}: una perturbación de demanda en un sector viaja \emph{aguas arriba}, hacia sus proveedores (anexo~\ref{anexo:redes}).
  \end{itemize}
\item \emph{La interacción con las perturbaciones de oferta}: las 2 rara vez llegan solas.
  \begin{itemize}
  \item \emph{De la oferta a la demanda}: si una perturbación cierra unos sectores, los bienes de los distintos sectores se sustituyen mal y los hogares no pueden asegurarse, la renta perdida reduce el gasto en los demás, y la caída de la demanda puede superar a la de la oferta. Es la \emph{perturbación de oferta keynesiana} (\emph{Keynesian supply shock}) de la Covid; con un solo sector no puede darse\nocite{GuerrieriEtAl2022} \vertema{3.A.40}.
  \item \emph{De la demanda a la oferta}: en 2021 la demanda se recuperó deprisa, empujada por la política fiscal y el ahorro acumulado, y se encontró con una oferta atascada: con la OA casi vertical en los sectores con cuellos de botella (\emph{bottlenecks}), la demanda se convirtió en precios. El origen de la inflación de 2021-2023 fueron la energía, los alimentos y los atascos; su persistencia, un mercado de trabajo muy tenso\nocite{BernankeBlanchard2025What}.
  \item Por eso no hay una política para cada perturbación «pura»: hay que saber qué parte de lo que se ve es de demanda y qué parte de oferta.
  \end{itemize}
\end{itemize}

\begin{anotaciones}
\textbf{La Covid: 2 respuestas a una perturbación de oferta keynesiana}
\begin{itemize}
\item La respuesta que pide no es un estímulo genérico, sino ayudas a los trabajadores sin liquidez, que gastan lo que reciben, y a las empresas para conservar los empleos\nocite{GuerrieriEtAl2022}.
\item Estados Unidos reforzó las prestaciones por desempleo y dio transferencias directas: su tasa de paro pasó del 3,5~\% en febrero de 2020 al 14,8~\% en abril\nocite{FRED2026UNRATE}.
\item La zona del euro optó sobre todo por mantener el vínculo laboral con regulaciones temporales de empleo, como los ERTE en España: su paro pasó del 7,3~\% al 7,4~\%, y el de España, del 13,6~\% al 15,3~\%\nocite{Eurostat2026unertm} \vertema{4.A.29}. Conservar el empleo conserva el capital humano propio de cada puesto para la recuperación.
\end{itemize}
\end{anotaciones}

% --------------------------------------------------------------------------
\subsection{La respuesta a las perturbaciones de demanda}\label{sec:demanda-respuesta}

\begin{itemize}
\item \emph{La política monetaria: la divina coincidencia}. Con una perturbación de demanda y sin perturbaciones de costes, basta con que el tipo de interés siga al natural: la DA vuelve a su sitio y la economía regresa a $E_0$, con la inflación en el objetivo y la brecha cerrada a la vez (gráfico~\ref{fig:nek-perturbaciones}, panel d). Un instrumento basta para 2 objetivos porque se mueven juntos. Es el fundamento de los objetivos de inflación \vertema{3.A.36}.
\item \emph{Cuando la monetaria no llega: la política fiscal}.
  \begin{itemize}
  \item \emph{Su multiplicador depende de la regla}: en la forma reducida \eqref{eq:forma-reducida}, el gasto sube $r^n_t$ y la brecha, pero tanto menos cuanto más agresiva es la respuesta del banco central a la inflación.
  \item \emph{En el límite inferior de los tipos}, el banco central no los sube tras la expansión, el gasto sube la inflación esperada y baja el tipo real, y el multiplicador puede ser bastante mayor que 1\nocite{ChristianoEichenbaumRebelo2011}: es el caso extremo de que la regla monetaria decide el multiplicador.
  \item \emph{Sin moneda propia}, la fiscal es el único instrumento para las perturbaciones de cada país (\ref{sec:externas}) \vertema{4.B.4}.
  \item \emph{Si los hogares son ricardianos}, una expansión financiada con deuda se compensa en parte con más ahorro; si una parte vive al día, el efecto es mayor\nocite{GaliLopezSalidoValles2007}.
  \end{itemize}
\item \emph{La combinación de políticas} (\emph{policy mix}): con el tono de la monetaria medido por el tipo real frente al natural y el de la fiscal por el impulso fiscal, hay 4 cuadrantes (figura~\ref{figu:policy-mix}) \vertema{3.A.40}.
  \begin{itemize}
  \item Ante una perturbación de demanda, las 2 deben ir en el mismo sentido, como en 2009 y en 2020; ante una de costes con expectativas en peligro, la monetaria contractiva y la fiscal solo focalizada, como en 2022-2023.
  \item Detrás está la pregunta de quién se ajusta para cumplir la restricción del Estado: con dominación monetaria, la fiscal; con dominación fiscal, la monetaria\nocite{SargentWallace1981Some,Leeper1991} \vertema{3.A.39}.
  \end{itemize}
\end{itemize}

\begin{infografia}[H]
\centering
\caption{La combinación de políticas (\emph{policy mix}): los 4 cuadrantes}\label{figu:policy-mix}
\grafico[0.6\textwidth]{policy-mix}
\fuente{Elaboración propia. Ejemplos de cada cuadrante: las 2 expansivas, en 2009 y 2020; las 2 contractivas, en 2022-2023; fiscal expansiva con monetaria contractiva, los Estados Unidos de los 80 y el Reino Unido en 2022; fiscal contractiva con monetaria expansiva, la zona del euro desde 2012. Que la combinación sea coherente no exige que las 2 políticas vayan en el mismo sentido.}
\end{infografia}

\begin{itemize}
\item \emph{Recapitulación de II}: el origen decide la política. Ante una perturbación de demanda (o tecnológica), la monetaria estabiliza a la vez renta e inflación, y si no puede, la fiscal; ante una de costes, la política de demanda solo reparte el golpe y la de oferta es la única que mejora las 2.
\item Hasta aquí, la economía era cerrada o el exterior estaba dado. Pero en una economía como la española muchas perturbaciones llegan de fuera, y la respuesta depende de algo que la economía cerrada no tiene: el régimen cambiario.
\end{itemize}


% ==========================================================================
\section{Extensión: las perturbaciones en una economía abierta}\label{sec:abierta}
\minutos{5}

\begin{itemize}
\item Abrimos el modelo neokeynesiano para ver cómo llegan las perturbaciones de fuera y qué régimen cambiario y qué regla de política las amortiguan mejor.
\item \emph{Conclusión anticipada}: la apertura cambia los coeficientes, no la forma; el banco central debe estabilizar la inflación de lo que se produce y dejar flotar el tipo de cambio, que amortigua las perturbaciones reales de demanda pero no las financieras ni, para un importador, las de la energía; sin moneda propia, la política fiscal y las de oferta hacen el trabajo.
\end{itemize}

% --------------------------------------------------------------------------
\subsection{El modelo neokeynesiano de una economía abierta}\label{sec:gali-monacelli}

\paragraph{a) Idea}

\begin{itemize}
\item Con \textcite{ObstfeldRogoff1995Redux} nace la \concepto{nueva macroeconomía abierta} (\emph{New Open Economy Macroeconomics}), que traslada a la economía abierta los modelos con agentes que optimizan, competencia monopolística y precios rígidos.
\item \textcite{GaliMonacelli2005Monetary} construyen la versión para una economía pequeña del modelo del apartado~\ref{sec:nek}: es la versión moderna de \autor{Mundell}-\autor{Fleming}. Responden a 2 preguntas: ¿qué cambia la apertura en la IS dinámica y en la curva de \autor{Phillips}? ¿Qué inflación debe estabilizar el banco central, la de lo que se produce o la de lo que se consume?
\end{itemize}

\paragraph{b) Modelo}

\subparagraph{i) Supuestos}

\begin{itemize}
\item \emph{Un continuo de economías pequeñas}: ninguna influye en las demás, y la producción mundial, $y^*_t$ (en logaritmos), está dada para cada una.
\item \emph{Hogares}: los del apartado~\ref{sec:nek}; consumen bienes nacionales y extranjeros, con una parte $\alpha$ de importaciones: $\alpha$ es el \concepto{grado de apertura} y $1-\alpha$, el sesgo hacia lo nacional.
\item \emph{Empresas}: las del apartado~\ref{sec:nek} (competencia monopolística y \autor{Calvo}).
\item \emph{Precios en la moneda del productor} (\emph{producer currency pricing}) y ley del precio único: una depreciación encarece en la misma proporción lo que se importa.
\item \emph{Mercados financieros completos}: los hogares de todos los países se aseguran entre sí.
\end{itemize}

\subparagraph{ii) Desarrollo}

\begin{itemize}
\item \emph{2 inflaciones}. Con $\tau_t\equiv p_{F,t}-p_{H,t}$ el precio de los bienes importados en términos de los nacionales, en logaritmos,\footnote{$\tau_t$ es la inversa de la relación real de intercambio ($P_X/P_M$): sube cuando esta empeora. Con la ley del precio único, $\tau_t=e_t+p^*_t-p_{H,t}$, con $e_t\equiv\ln E_t$, y la identidad \eqref{eq:inflacion-gm} es la misma que \eqref{eq:inflacion-abierta}.} la inflación de consumo es
  \begin{equation}
  \pi_t=\pi_{H,t}+\alpha\,\Delta\tau_t . \label{eq:inflacion-gm}
  \end{equation}
\item \emph{La IS dinámica abierta}: de la ecuación de \autor{Euler}, el reparto del riesgo entre países y el vaciado del mercado de bienes nacionales (anexo~\ref{anexo:nek}),
  \funcionclave[eq:is-abierta-gm]{x_t=\mathbb{E}_t[x_{t+1}]-\frac{1}{\theta_{\alpha}}\Big(i_t-\mathbb{E}_t[\pi_{H,t+1}]-r^n_t\Big),\qquad \theta_{\alpha}\equiv\frac{\theta}{1+\alpha\,\Theta}}
  con $\Theta$ un término que depende de la apertura y de las elasticidades de sustitución entre bienes nacionales y extranjeros. Si son altas ($\Theta>0$), la apertura hace la brecha más sensible al tipo de interés, porque una subida de tipos también aprecia la moneda y resta exportaciones. El tipo natural depende ahora también del crecimiento esperado de la producción mundial.
\item \emph{La curva de \autor{Phillips} abierta}, para la inflación de los bienes nacionales, que es la que fijan las empresas:
  \funcionclave[eq:nkpc-abierta]{\pi_{H,t}=\beta\,\mathbb{E}_t[\pi_{H,t+1}]+\kappa_{\alpha}\,x_t+u_t,\qquad \kappa_{\alpha}\equiv\frac{(1-\eta)(1-\beta\eta)}{\eta}\,(\theta_{\alpha}+\varphi)}
\item \emph{La versión gráfica} es la del apartado~\ref{sec:nek} en el plano $(x_t,\pi_{H,t})$, con $\theta_\alpha$ y $\kappa_\alpha$ (gráfico~\ref{fig:nek-perturbaciones}, con los mismos supuestos), y la forma reducida \eqref{eq:forma-reducida} con esos coeficientes.
\end{itemize}

\subparagraph{iii) Implicaciones}

\begin{itemize}
\item \emph{Misma estructura que la economía cerrada}: la apertura cambia los coeficientes ($\theta_{\alpha}$, $\kappa_{\alpha}$) y el tipo natural, no la forma. Con $\alpha=0$, o con elasticidades unitarias, es el modelo canónico\nocite{Gali2015Monetary}.
\item \emph{El tipo de cambio es un precio relativo}: con precios rígidos, mueve la relación de intercambio y reasigna el gasto entre bienes nacionales y extranjeros, como en \autor{Mundell}-\autor{Fleming}, pero ahora sale de la optimización. Una contracción monetaria sube el tipo real, aprecia la moneda y mejora la relación real de intercambio\nocite{Gali2015Monetary}.
\item \emph{2 inflaciones, 2 objetivos posibles}: la divina coincidencia vale para $\pi_{H,t}$, no para la inflación de consumo, que incluye la relación de intercambio. Elegir entre ellas es la primera decisión del banco central de una economía abierta.
\end{itemize}

\paragraph{c) Extensiones}

\begin{itemize}
\item \emph{Precios en la moneda del comprador} (\emph{local currency pricing}): los precios de importación no reaccionan a corto plazo al tipo de cambio, y el banco central ya no puede cerrar a la vez la brecha y el desalineamiento de precios\nocite{CorsettiPesenti2009} \vertema{3.A.37} \vertema{3.B.19}. Con la \emph{fijación de precios en la moneda dominante} (\emph{dominant currency pricing}), los tipos de la Reserva Federal mueven los precios de importación en países que ni siquiera comercian con Estados Unidos\nocite{Gopinath2020DCP}.
\item \emph{Mercados incompletos y fricciones financieras}: sin seguro entre países, la deuda exterior importa, y en los emergentes el tipo de interés mundial y las primas de riesgo dominan el ciclo\nocite{UribeSchmittGrohe2017Open}.
\item \emph{Uniones monetarias}: \textcite{GaliMonacelli2008Optimal} \vertema{3.B.16}.
\end{itemize}

\paragraph{d) Evidencia empírica}

\begin{itemize}
\item \emph{El reparto del riesgo falla}: con mercados completos, el país cuya moneda se deprecia en términos reales debería consumir relativamente menos; en los datos, la correlación entre el consumo relativo y el tipo de cambio real es nula o negativa\nocite{BackusSmith1993Consumption}.
\item \emph{El traspaso es incompleto}: la mayor parte del comercio mundial se factura en unas pocas monedas, sobre todo el dólar, y los precios de importación en moneda local apenas se mueven a corto plazo\nocite{Gopinath2020DCP}.
\end{itemize}

\paragraph{e) Valoración}

\begin{itemize}
\item \emph{Lo que aporta}: es el marco con el que se piensa hoy la política monetaria de una economía abierta, con un criterio de bienestar para comparar reglas y regímenes.
\item \emph{Sus límites}: mercados completos, traspaso completo y economía pequeña; sin deuda exterior ni efectos sobre los demás países. Con él, vemos las perturbaciones que llegan de fuera.
\end{itemize}

% --------------------------------------------------------------------------
\subsection{Las perturbaciones externas y el régimen cambiario}\label{sec:externas}

\begin{itemize}
\item \emph{Las perturbaciones que llegan de fuera}, y cómo las amortigua el régimen:
  \begin{itemize}
  \item \emph{una caída de la demanda exterior} es de demanda: con tipo flexible, la depreciación la compensa (el aislamiento de \autor{Mundell}-\autor{Fleming}; en el modelo neokeynesiano, el banco central baja los tipos siguiendo al natural, que cae con $y^*_t$); con tipo fijo, cae la renta y queda la política fiscal;
  \item \emph{una subida del tipo de interés mundial} es financiera: con tipo fijo sube el tipo propio; con flexible, deprecia la moneda, pero por el ciclo financiero global (\emph{global financial cycle}) contrae el crédito y el precio de los activos. Más que un trilema hay un dilema: la elección real sería entre la movilidad del capital y la autonomía monetaria\nocite{Rey2013Dilema,MirandaAgrippinoRey2020} \vertema{3.A.37}, y un país cuya moneda no es la del centro solo conserva su autonomía si gestiona su cuenta financiera, directamente o con la política macroprudencial;
  \item \emph{un encarecimiento de la energía} es de oferta y llega por el precio en dólares y por el tipo de cambio \eqref{eq:inflacion-abierta}. Para un importador, la flotación puede amplificarlo: con mercados incompletos, el encarecimiento transfiere riqueza a los exportadores, reduce el consumo del importador y deprecia su tipo real, y la energía se encarece todavía más en moneda nacional; para un exportador, la apreciación lo amortigua\nocite{BodensteinErcegGuerrieri2011}.
  \end{itemize}
\item \emph{El régimen como amortiguador}: con tipo flexible, una perturbación real que exige cambiar los precios relativos se ajusta por el tipo de cambio, sin esperar a los salarios\nocite{Friedman1953Case}; es la lógica de que más vale doblarse que quebrarse. Con tipo fijo, el ajuste es interno y lento (\ref{sec:da-oa}). La tabla~\ref{tab:perturbaciones} del anexo~\ref{anexo:perturbaciones} resume la respuesta a cada perturbación según el régimen.
\item \emph{¿Qué inflación debe mirar el banco central?} \textcite{GaliMonacelli2005Monetary} comparan 3 reglas:
  \begin{equation}
  \text{DITR:}\ \ i_t=\rho+\phi_{\pi}\,\pi_{H,t};\qquad \text{CITR:}\ \ i_t=\rho+\phi_{\pi}\,\pi_t;\qquad \text{PEG:}\ \ e_t=0, \label{eq:reglas-gm}
  \end{equation}
  una regla de \autor{Taylor} con la inflación interna (\emph{domestic inflation \autor{Taylor} rule}), otra con la de consumo (\emph{CPI inflation \autor{Taylor} rule}) y un tipo de cambio fijo (\emph{peg}) \vertema{3.A.36}.
  \begin{itemize}
  \item \emph{La política óptima}, en el caso en que pueden calcular el bienestar (elasticidades unitarias), es estabilizar del todo la inflación interna: $\pi_{H,t}=0$ y $x_t=0$. El tipo de cambio y la relación de intercambio fluctúan mucho, y eso es bueno: hacen el ajuste de los precios relativos que los precios rígidos no hacen.
  \item \emph{La ordenación}: la regla con la inflación interna es mejor que la de la inflación de consumo, y esta, mejor que el tipo fijo. Cuanto más suaviza una regla el tipo de cambio, más volátiles son la brecha y la inflación interna, y mayor la pérdida de bienestar.
  \item \emph{Volvamos a la energía}: la regla óptima mira la inflación de lo que se produce dentro y deja pasar el efecto directo de lo importado. Es el fundamento, en una economía abierta, de mirar a través de una perturbación energética mientras no contagie a los precios y salarios nacionales (\ref{sec:oferta-respuesta}).
  \item \emph{Pero el miedo a flotar} (\emph{fear of floating}) \emph{tiene razones}: si los precios de importación se fijan en la moneda del comprador o en dólares, el tipo de cambio no reasigna el gasto a corto plazo, y la política óptima debe atender también a sus desalineamientos\nocite{Engel2011Currency}. Muchos países con tipo flexible intervienen para suavizarlo\nocite{CalvoReinhart2002}.
  \end{itemize}
\item \emph{Sin moneda propia: la unión monetaria}. El banco central común responde a la media de la unión, y frente a una perturbación propia el instrumento de cada país es la política fiscal.
  \begin{itemize}
  \item En un modelo neokeynesiano de unión monetaria, lo óptimo es que el banco central estabilice la inflación media y que cada gobierno use el gasto para cerrar su brecha\nocite{GaliMonacelli2008Optimal} \vertema{3.B.16}. Los datos lo apoyan: el multiplicador \emph{relativo} del gasto de una región dentro de una unión monetaria está en torno a 1,5\nocite{NakamuraSteinsson2014,FarhiWerning2016}.
  \item La apertura reduce el multiplicador, porque parte del gasto se va en importaciones; en los datos es mayor con tipo predeterminado que con tipo flexible\nocite{Furlanetto2006Fiscal,IlzetzkiMendozaVegh2013}, aunque con tipo fijo, si el gasto sube la inflación, el tipo real a largo plazo sube y frena la demanda privada\nocite{CorsettiKuesterMuller2011}.
  \item \emph{La devaluación fiscal} (\emph{fiscal devaluation}): subir el IVA y bajar las cotizaciones sociales de las empresas abarata las exportaciones y encarece las importaciones, como una devaluación nominal, sin tocar el tipo de cambio\nocite{FarhiGopinathItskhoki2014}. Es, con las políticas de oferta, la alternativa a la devaluación interna.
  \end{itemize}
\item \emph{El episodio de 2021-2023 y 2026 en la zona del euro}:
  \begin{itemize}
  \item La depreciación del euro amplificó el encarecimiento de la energía (gráfico~\ref{fig:energia-tipo-cambio}): entre enero de 2021 y junio de 2022 el Brent subió un 124~\% en dólares y un 158~\% en euros. La inflación tocó techo en el 10,6~\% en octubre de 2022 y la subyacente, con retraso, en el 5,7~\% en marzo de 2023\nocite{Eurostat2026IAPC}.
  \item La Reserva Federal empezó a subir los tipos el 16 de marzo de 2022\nocite{Fed2022FOMCMarzo}, y el BCE, por primera vez en 11 años, el 21 de julio de 2022 y, en 10 subidas, hasta septiembre de 2023, llevó la facilidad de depósito del $-0{,}50$~\% al 4~\%\nocite{BCE2022DecisionesJulio,BCE2026TiposOficiales}. La diferencia tiene base: en la zona del euro pesan más los efectos de segunda ronda del petróleo sobre los salarios, y el BCE da más peso a la inflación frente a la producción\nocite{PeersmanVanRobays2009}.
  \item En 2026, con la lección de 2022, el BCE ha subido la facilidad de depósito al 2,25~\% en junio y al 2,50~\% en septiembre, aunque la subyacente no se ha movido: en octubre la energía se encarece un 20,8~\% en un año, la inflación está en el 4,0~\% y la subyacente en el 2,5~\%\nocite{BCE2026TiposOficiales,Eurostat2026IAPCoct}. Con el euro más fuerte, el Brent en euros ha subido menos que en dólares.
  \end{itemize}
\end{itemize}

\begin{anotaciones}
\textbf{Las cifras de 2022}
\begin{itemize}
\item El gas europeo costó más de 8 veces lo que en 2019, y en agosto llegó a unos 236 euros por megavatio hora\nocite{BancoMundial2026PinkSheet}; la energía se encareció un 44,3~\% interanual en marzo en la zona del euro\nocite{Eurostat2026IAPC}. La zona del euro pasó de crecer un 3,7~\% en 2022 a un 0,4~\% en 2023, y España, del 6,4~\% al 2,4~\%\nocite{Eurostat2026nama10gdp,INE2026CNA2025}.
\item \emph{La amplificación por el tipo de cambio}: el Brent subió más en euros que en dólares\nocite{EIA2026Brent} porque, entre enero de 2021 y junio de 2022, el euro se depreció un 13~\% frente al dólar, y siguió cayendo hasta 0,98 dólares en octubre de 2022, un 19~\% menos que en enero de 2021\nocite{BCE2026EXR}.
\end{itemize}
\begin{figuracaja}
\centering
\caption{El petróleo en dólares y en euros y el tipo de cambio del euro, 2019-2026}\label{fig:energia-tipo-cambio}
\grafico[0.92\textwidth]{energia-tipo-cambio-zona-euro}
\fuente{Elaboración propia con datos de \textcite{EIA2026Brent} (Brent, dólares por barril) y \textcite{BCE2026EXR} (dólares por euro, media mensual); el Brent en euros es el precio en dólares dividido por el tipo de cambio. Medias mensuales, enero de 2019-octubre de 2026. En 2022 las 2 líneas del Brent se separan: el euro se deprecia y la energía sube más en euros que en dólares.}
\end{figuracaja}
\end{anotaciones}

\begin{itemize}
\item \emph{Recapitulación de III}: en una economía abierta la estructura del modelo es la misma, pero las perturbaciones llegan también por el tipo de cambio y el tipo de interés mundial. La flotación amortigua las reales de demanda; el banco central debe mirar la inflación interna, y sin moneda propia la política fiscal, la devaluación fiscal y las políticas de oferta sustituyen al tipo de cambio.
\end{itemize}

% ==========================================================================
\section*{Conclusión}
\minutos{1}

\begin{itemize}
\item \textbf{\emph{Recapitulación}}:
  \begin{itemize}
  \item En conclusión, la renta y la inflación de una economía abierta se determinan con una demanda agregada que nace de \autor{Mundell}-\autor{Fleming} y depende del tipo de cambio real, y una oferta que paga lo importado; su versión moderna es el modelo neokeynesiano, que distingue las perturbaciones de demanda de las de costes.
  \item Ante una perturbación de demanda, la política monetaria (o, sin ella, la fiscal) estabiliza a la vez renta e inflación; ante una de costes, la política de demanda solo reparte el golpe, y las de oferta son las únicas que mejoran las 2.
  \item En una economía abierta, el régimen cambiario decide qué instrumento queda y qué inflación conviene estabilizar.
  \end{itemize}
\item \textbf{\emph{Relevancia}}:
  \begin{itemize}
  \item Es el marco con el que el BCE y los gobiernos de la zona del euro han respondido a la Covid y a la energía cara de 2022 y de 2026: una política monetaria común, una política fiscal nacional y unas reformas que deciden la competitividad.
  \end{itemize}
\item \textbf{\emph{Extensiones}}:
  \begin{itemize}
  \item La determinación del tipo de cambio \vertema{3.B.14}, la elección del régimen cambiario y las crisis \vertema{3.B.15} \vertema{3.B.17}, las uniones monetarias \vertema{3.B.16} y la coordinación internacional de las políticas \vertema{3.B.19}.
  \item La economía española dentro del euro y su competitividad \vertema{4.A.22}.
  \end{itemize}
\item \textbf{\emph{Opinión}}:
  \begin{itemize}
  \item \autor{Mundell}-\autor{Fleming} sigue siendo la mejor primera aproximación, y el modelo neokeynesiano la mejor guía para la política; pero las perturbaciones de verdad llegan mezcladas, y saber qué parte es de oferta y qué parte de demanda es la mitad de la respuesta.
  \end{itemize}
\item \textbf{\emph{Idea final}}:
  \begin{itemize}
  \item Las desgracias nunca vienen solas: en 2020-2022 llegaron juntas una de oferta, una de demanda y otra de fuera, por el precio de lo que se importa y por el tipo de cambio con que se paga.
  \item Las políticas de demanda deciden cómo se reparte el golpe, y su eficacia depende del régimen cambiario; solo las de oferta hacen que la próxima desgracia duela menos.
  \end{itemize}
\end{itemize}

\clearpage
\appendix

% ==========================================================================
\section{Anexo: el modelo de \autor{Mundell}-\autor{Fleming} paso a paso}\label{anexo:mundell-fleming}

\begin{itemize}
\item \emph{Las 3 ecuaciones}, con $P=P^*=1$ (luego $q=E$):
  \begin{align*}
  \text{IS:}&\quad \big[1-c\,(1-t)+m\big]\,Y=\overline{A}-b\,i+x\,Y^*+\nu\,E,\\
  \text{LM:}&\quad M=k\,Y-h\,i,\\
  \text{BP:}&\quad x\,Y^*-m\,Y+\nu\,E+\beta_K\,(i-i^*)=0 .
  \end{align*}
\item \emph{Tipo fijo} ($E=\overline{E}$; incógnitas $Y$, $i$ y $M$). De la BP, $i=i^*+(m\,Y-x\,Y^*-\nu\,\overline{E})/\beta_K$; llevándolo a la IS y multiplicando por $\beta_K$,
  \begin{equation*}
  Y=\frac{\beta_K\,\overline{A}-b\,\beta_K\,i^*+(b+\beta_K)\,(x\,Y^*+\nu\,\overline{E})}{D_F},\qquad D_F\equiv\big[1-c\,(1-t)\big]\,\beta_K+m\,(b+\beta_K),
  \end{equation*}
  con $D_F$ local de este anexo, y la cantidad de dinero sale de la LM. Los multiplicadores son
  \begin{equation*}
  \frac{\partial Y}{\partial\overline{G}}=\frac{\beta_K}{D_F},\quad \frac{\partial Y}{\partial\overline{E}}=\frac{\nu\,(b+\beta_K)}{D_F},\quad \frac{\partial Y}{\partial Y^*}=\frac{x\,(b+\beta_K)}{D_F},\quad \frac{\partial Y}{\partial i^*}=-\frac{b\,\beta_K}{D_F},\quad \frac{\partial Y}{\partial M}=0 .
  \end{equation*}
  \begin{itemize}
  \item La compra de bonos por el banco central sube el crédito interno y baja las reservas en lo mismo: $M$ no cambia.
  \item $\partial(\beta_K/D_F)/\partial\beta_K=m\,b/D_F^2>0$: la fiscal gana eficacia con la movilidad.
  \item Con $\beta_K\to\infty$, $\partial Y/\partial\overline{G}\to1/[1-c\,(1-t)+m]$, el multiplicador abierto con $i=i^*$. Con $\beta_K=0$, $\partial Y/\partial\overline{G}=0$ y $Y=(x\,Y^*+\nu\,\overline{E})/m$, la renta que equilibra el comercio.
  \end{itemize}
\item \emph{Tipo flexible} ($M$ dado; incógnitas $Y$, $i$ y $E$). La BP da $x\,Y^*+\nu\,E=m\,Y-\beta_K\,(i-i^*)$; llevándolo a la IS, se cancelan $m\,Y$ y el comercio:
  \begin{equation*}
  \big[1-c\,(1-t)\big]\,Y=\overline{A}-(b+\beta_K)\,i+\beta_K\,i^* ,
  \end{equation*}
  una IS de economía cerrada con sensibilidad $b+\beta_K$. Con la LM, $i=(k\,Y-M)/h$, y multiplicando por $h$,
  \begin{equation*}
  Y=\frac{h\,\overline{A}+(b+\beta_K)\,M+h\,\beta_K\,i^*}{D_X},\qquad D_X\equiv h\,\big[1-c\,(1-t)\big]+(b+\beta_K)\,k ,
  \end{equation*}
  y el tipo de cambio sale de la BP, $\nu\,E=m\,Y-x\,Y^*-\beta_K\,(i-i^*)$. Los multiplicadores son
  \begin{equation*}
  \frac{\partial Y}{\partial\overline{G}}=\frac{h}{D_X},\quad \frac{\partial Y}{\partial M}=\frac{b+\beta_K}{D_X},\quad \frac{\partial Y}{\partial Y^*}=0,\quad \frac{\partial Y}{\partial i^*}=\frac{h\,\beta_K}{D_X}>0 .
  \end{equation*}
  \begin{itemize}
  \item $\partial(h/D_X)/\partial\beta_K=-h\,k/D_X^2<0$ y $\partial[(b+\beta_K)/D_X]/\partial\beta_K=h\,[1-c\,(1-t)]/D_X^2>0$.
  \item Con $\beta_K=0$, los multiplicadores son los del IS-LM cerrado sin importaciones, $1/[1-c\,(1-t)+b\,k/h]$ y $b/\{h\,[1-c\,(1-t)]+b\,k\}$ \vertema{3.A.5}: el tipo de cambio mantiene el comercio equilibrado y anula la fuga de las importaciones. Con $\beta_K\to\infty$, la fiscal da 0 y la monetaria $1/k$.
  \item \emph{Aislamiento}: $Y$ no depende de $Y^*$, porque una caída de la demanda exterior deprecia la moneda lo justo para compensarla. Una subida de $i^*$ sí sube la renta en el modelo, porque deprecia la moneda; en los datos, en cambio, la subida de los tipos de la Reserva Federal contrae a los demás por el canal financiero\nocite{MirandaAgrippinoRey2020}.
  \end{itemize}
\item \emph{Los déficits gemelos}. Con tipo flexible y $\beta_K\to\infty$, $i=i^*$ e $Y=(M+h\,i^*)/k$ no dependen de $\overline{G}$; como el consumo y la inversión tampoco cambian, $\Delta XN=\nu\,\Delta E=-\Delta\overline{G}$: la moneda se aprecia y el saldo comercial empeora en lo que sube el gasto. Con tipo fijo, $\Delta XN=-m\,\Delta Y=-m\,\Delta\overline{G}/[1-c\,(1-t)+m]$.
\end{itemize}
\findemo

% ==========================================================================
\section{Anexo: oferta y demanda agregadas abiertas paso a paso}\label{anexo:oa-da}

\begin{itemize}
\item \emph{La DA con movilidad perfecta}. Con $P$ variable, $q=EP^*/P$, $i=i^*$, y \eqref{eq:da-abierta}:
  \begin{itemize}
  \item con tipo fijo, $\dfrac{\partial Y}{\partial P}=-\dfrac{\nu\,\overline{E}P^*}{[1-c\,(1-t)+m]\,P^2}<0$ (solo el efecto competitividad);
  \item con tipo flexible, $\dfrac{\partial Y}{\partial P}=-\dfrac{M}{k\,P^2}<0$ (solo el efecto saldos reales).
  \end{itemize}
  Con movilidad imperfecta actúan los 2 canales.
\item \emph{El medio plazo con tipo fijo}. En $Y=Y^n$, la IS da $\nu\,\overline{E}P^*/P=[1-c\,(1-t)+m]\,Y^n-\overline{A}+b\,i^*-x\,Y^*$: más gasto exige un $q$ menor (una apreciación real), y $\nu\,\Delta q=-\Delta\overline{G}$.
\item \emph{El camino hacia el medio plazo con tipo fijo}, en 3 pasos\nocite{Blanchard2021Macroeconomics}:
  \begin{itemize}
  \item \emph{Paso 1}: con la inflación esperada igual a la exterior, $\bar\pi=\gamma_{P^*}$, la curva de \autor{Phillips} en producción es $\pi_t-\bar\pi=\kappa_Y\,(Y_t-Y^n)$, con $\kappa_Y>0$ (local).
  \item \emph{Paso 2}: con $E$ fijo, el tipo real cambia lo que la inflación exterior supera a la nacional: $\gamma_{q,t}\simeq\gamma_{P^*}-\pi_t=-\kappa_Y\,(Y_t-Y^n)$.
  \item \emph{Paso 3}: si $Y_t<Y^n$, $q$ sube (depreciación real), las exportaciones netas crecen y la renta se acerca a $Y^n$; el proceso solo se para con $Y_t=Y^n$ y $\pi_t=\gamma_{P^*}$. Una devaluación de la cuantía adecuada da ese mismo $q$ de una vez.
  \end{itemize}
\item \emph{La cuña exterior}. Con una parte $\alpha$ de importaciones en el consumo, el índice de precios de consumo es $P_C=P^{1-\alpha}(EP^*)^{\alpha}$ ($P_C$ local), luego $P_C/P=(EP^*/P)^{\alpha}=q^{\alpha}$. Con la curva de salarios escrita con el índice de consumo y la de precios del modelo WS-PS \vertema{3.A.28},
  \begin{equation*}
  \frac{W}{P_C^e}=F(U,Z)\cdot\mathrm{PAT}^e,\qquad \frac{W}{P}=\frac{\mathrm{PAT}}{1+\mu}\quad\Longrightarrow\quad P=(1+\mu)\,P_C^e\,F(U,Z)\,\frac{\mathrm{PAT}^e}{\mathrm{PAT}} .
  \end{equation*}
  \begin{itemize}
  \item \emph{A corto plazo}, con $P_C^e$ dado y $U$ que baja cuando sube $Y$, es la OA creciente; sube si suben los precios de consumo esperados (una depreciación esperada), el margen (la energía) o $Z$.
  \item \emph{A medio plazo}, con $P_C^e=P_C$ y $\mathrm{PAT}^e=\mathrm{PAT}$, $1=(1+\mu)\,q^{\alpha}F(U^n,Z)$, que es \eqref{eq:nairu-abierta}. Como $F_U<0$, un $q$ o un $\mu$ mayores bajan $F(U^n,Z)$ y suben $U^n$. Con $\alpha=0$ es la NAIRU de la economía cerrada.
  \item \emph{Un matiz}: con $E$ fijo, un $P$ mayor baja $q$ y la NAIRU, y la OA de largo plazo se inclina algo hacia la derecha; con $E$ flexible, $q$ lo fijan factores reales y es vertical. El cuerpo la dibuja vertical, para una cuña dada.
  \end{itemize}
\item \emph{La inflación de consumo}. En logaritmos, $\ln P_C=(1-\alpha)\ln P+\alpha\,(\ln E+\ln P^*)$; en primeras diferencias, con $\pi_{H,t}\equiv\Delta\ln P_t$, es \eqref{eq:inflacion-abierta}.
\item \emph{El modelo de demandas rivales} en el plano $(Y,q)$\nocite{LayardNickellJackman1991Unemployment,CarlinSoskice2015Macroeconomics} (gráfico~\ref{fig:demandas-rivales}):
  \begin{itemize}
  \item \emph{la EDR} (equilibrio de las demandas rivales) es \eqref{eq:nairu-abierta} escrita en producción: los pares $(Y,q)$ en los que trabajadores, empresas y proveedores extranjeros aceptan su parte de la renta y la inflación no se mueve. Es \emph{decreciente}: con más $q$, el salario de consumo que deja el margen es menor, y para que los salarios pedidos quepan hace falta más paro, menos $Y$. Por encima hay presiones inflacionistas; por debajo, deflacionistas;
  \item \emph{la BC} (equilibrio comercial) es creciente: más renta, más importaciones, y hace falta más competitividad;
  \item \emph{la DA} (equilibrio del mercado de bienes) es creciente y más inclinada que la BC: una depreciación sube la renta, pero, como hay otras fugas además de las importaciones, el comercio queda con superávit.
  \end{itemize}
  \emph{Un encarecimiento de la energía importada} desplaza las 3 a la izquierda: la EDR, porque baja el salario de consumo compatible con el margen; la BC, porque hace falta más depreciación para cada renta; la DA, porque parte del gasto se va en la energía. El nuevo equilibrio, $E'$, tiene menos producción y un tipo real más depreciado. Si se intenta compensar con estímulos, la inflación sigue subiendo (el error de 1973). Con tipo fijo, la inflación aprecia el tipo real y lleva la economía a la EDR con déficit comercial, y corregirlo exige ajuste fiscal y devaluación; con tipo flexible, la espiral de inflación y depreciación solo aparece si el banco central acomoda.
\end{itemize}

\begin{figure}[H]
\centering
\caption{El modelo de demandas rivales en una economía abierta ante un encarecimiento de la energía importada}\label{fig:demandas-rivales}
\grafico[0.6\textwidth]{demandas-rivales}
\fuente{Elaboración propia a partir de \textcite{LayardNickellJackman1991Unemployment} y \textcite{CarlinSoskice2015Macroeconomics}. EDR: equilibrio de las demandas rivales sobre la renta (inflación estable); BC: equilibrio de la balanza comercial; DA: equilibrio del mercado de bienes. El encarecimiento de la energía desplaza las 3 a la izquierda (en rojo) y el equilibrio pasa de $E$ a $E'$.}
\end{figure}
\findemo

% ==========================================================================
\section{Anexo: el modelo neokeynesiano paso a paso}\label{anexo:nek}

\begin{itemize}
\item \emph{La DA, la OA y la forma reducida en la economía cerrada}, en 3 pasos:
  \begin{itemize}
  \item \emph{Paso 1}: con expectativas dadas ($\mathbb{E}_t[x_{t+1}]=\mathbb{E}_t[\pi_{t+1}]=0$, perturbaciones transitorias), la IS es $x_t=-\frac{1}{\theta}(i_t-r^n_t)$ y la NKPC, $\pi_t=\kappa\,x_t+u_t$.
  \item \emph{Paso 2}: con la regla \eqref{eq:taylor} dentro de la IS, $(\theta+\phi_x)\,x_t=-\phi_{\pi}\,\pi_t+\hat r^n_t-\varepsilon_t$, con $\hat r^n_t\equiv r^n_t-\rho$: es la DA \eqref{eq:da-nek}, decreciente con pendiente $-(\theta+\phi_x)/\phi_{\pi}$.
  \item \emph{Paso 3}: sustituyendo la NKPC, $(\theta+\phi_x+\phi_{\pi}\kappa)\,x_t=\hat r^n_t-\phi_{\pi}\,u_t-\varepsilon_t$, que es \eqref{eq:forma-reducida}. Con $u_t=0$ e $i_t$ que sigue a $r^n_t$, $x_t=\pi_t=0$: la divina coincidencia.
  \end{itemize}
\item \emph{La perturbación tecnológica}: con precios flexibles, la producción natural es $y^n_t=\psi\,a_t$, con $a_t\equiv\ln A_t$ y $\psi\equiv(1+\varphi)/(\theta+\varphi)$ (local), y el tipo natural, $r^n_t=\rho+\theta\,\mathbb{E}_t[\Delta y^n_{t+1}]$ más los términos de demanda. Con una mejora transitoria, $\mathbb{E}_t[\Delta y^n_{t+1}]<0$ y $r^n_t$ baja: la DA se desplaza a la izquierda.
\item \emph{La economía abierta} (\autor{Galí} y \autor{Monacelli}). Notación: minúsculas en logaritmos; $\sigma_{HF}$ es la elasticidad de sustitución entre bienes nacionales y extranjeros y $\sigma_F$ la que hay entre bienes de distintos países extranjeros (locales).
  \begin{itemize}
  \item \emph{El consumo y sus precios}:
  \begin{equation*}
  C_t=\Big[(1-\alpha)^{\frac{1}{\sigma_{HF}}}C_{H,t}^{\frac{\sigma_{HF}-1}{\sigma_{HF}}}+\alpha^{\frac{1}{\sigma_{HF}}}C_{F,t}^{\frac{\sigma_{HF}-1}{\sigma_{HF}}}\Big]^{\frac{\sigma_{HF}}{\sigma_{HF}-1}},\qquad P_t=\Big[(1-\alpha)P_{H,t}^{1-\sigma_{HF}}+\alpha\,P_{F,t}^{1-\sigma_{HF}}\Big]^{\frac{1}{1-\sigma_{HF}}}.
  \end{equation*}
  Log-linealizando, $p_t=(1-\alpha)\,p_{H,t}+\alpha\,p_{F,t}=p_{H,t}+\alpha\,\tau_t$, y en diferencias, \eqref{eq:inflacion-gm}.
  \item \emph{Ley del precio único y tipo de cambio real}: $p_{F,t}=e_t+p^*_t$, luego $\tau_t=e_t+p^*_t-p_{H,t}$, que es el logaritmo del $q$ del apartado~\ref{sec:da-oa} (con el deflactor). El tipo real con los índices de consumo es $e_t+p^*_t-p_t=(1-\alpha)\,\tau_t$.
  \item \emph{Reparto del riesgo}: con mercados completos, $c_t=c^*_t+\frac{1}{\theta}(1-\alpha)\,\tau_t$, y como el mundo es cerrado, $c^*_t=y^*_t$.
  \item \emph{Vaciado del mercado de bienes nacionales}: $y_t=c_t+\frac{\alpha\,\Psi}{\theta}\,\tau_t$, con $\Psi\equiv\theta\sigma_F+(1-\alpha)(\theta\sigma_{HF}-1)$ (local). Juntando las 2,
  \begin{equation*}
  y_t=y^*_t+\frac{1}{\theta_{\alpha}}\,\tau_t,\qquad \theta_{\alpha}\equiv\frac{\theta}{(1-\alpha)+\alpha\,\Psi}=\frac{\theta}{1+\alpha\,\Theta},\qquad \Theta\equiv\Psi-1=(\theta\sigma_F-1)+(1-\alpha)(\theta\sigma_{HF}-1).
  \end{equation*}
  \item \emph{La IS}: la ecuación de \autor{Euler}, $c_t=\mathbb{E}_t[c_{t+1}]-\frac1\theta(i_t-\mathbb{E}_t[\pi_{t+1}]-\rho)$, con las 2 relaciones anteriores y \eqref{eq:inflacion-gm}, da
  \begin{equation*}
  y_t=\mathbb{E}_t[y_{t+1}]-\frac{1}{\theta_{\alpha}}\big(i_t-\mathbb{E}_t[\pi_{H,t+1}]-\rho\big)+\alpha\,\Theta\,\mathbb{E}_t[\Delta y^*_{t+1}] .
  \end{equation*}
  Restando la misma ecuación con precios flexibles ($y^n_t$) sale \eqref{eq:is-abierta-gm}, con $r^n_t=\rho+\theta_{\alpha}\,\mathbb{E}_t[\Delta y^n_{t+1}]+\alpha\,\theta_{\alpha}\Theta\,\mathbb{E}_t[\Delta y^*_{t+1}]$.
  \item \emph{La curva de \autor{Phillips}}: el coste marginal real depende de la producción propia y de la mundial, y en desviación de la economía con precios flexibles es $\widehat{\mathrm{cmg}}_t=(\theta_{\alpha}+\varphi)\,x_t$. Con \autor{Calvo}, $\pi_{H,t}=\beta\,\mathbb{E}_t[\pi_{H,t+1}]+\frac{(1-\eta)(1-\beta\eta)}{\eta}\widehat{\mathrm{cmg}}_t$, que es \eqref{eq:nkpc-abierta} (sin $u_t$).
  \item \emph{El caso especial}: con $\theta=\sigma_{HF}=\sigma_F=1$, se tiene $\Psi=1$, $\Theta=0$, $\theta_{\alpha}=\theta$ y $\kappa_{\alpha}=\kappa$: el modelo es el cerrado. Con un subsidio al empleo que corrige el margen y el incentivo a mejorar la relación de intercambio, la asignación con precios flexibles es eficiente y la política óptima es $\pi_{H,t}=0$, $x_t=0$ en todo periodo\nocite{GaliMonacelli2005Monetary}.
  \item \emph{Las reglas} \eqref{eq:reglas-gm}: con la regla de la inflación interna, el tipo de cambio y $\tau_t$ fluctúan menos que con la óptima; con la del consumo, menos todavía, y con el tipo fijo el nominal no se mueve. Cuanto menos fluctúa $\tau_t$, más fluctúan la brecha y la inflación interna, y mayor es la pérdida: el tipo fijo es, en todas las calibraciones, el peor\nocite{GaliMonacelli2005Monetary}.
  \end{itemize}
\end{itemize}
\findemo

% ==========================================================================
\section{Anexo: perturbaciones sectoriales y redes de producción}\label{anexo:redes}

\begin{itemize}
\item \emph{El modelo}\nocite{AcemogluAkcigitKerr2016}: $n$ industrias producen en competencia perfecta con una tecnología \autor{Cobb}-\autor{Douglas}, $y_i=e^{z_i}\,l_i^{\alpha^l_i}\prod_j x_{ij}^{a_{ij}}$, con $\alpha^l_i+\sum_j a_{ij}=1$, donde $x_{ij}$ es lo que la industria $i$ compra a la $j$ y $z_i$ una perturbación de productividad (locales de este anexo). De la maximización del beneficio, $a_{ij}=p_j\,x_{ij}/(p_i\,y_i)$: la matriz $\mathbf A=(a_{ij})$ es la de \emph{insumo-producto}, y su inversa de \autor{Leontief}, $\mathbf H\equiv(\mathbf I-\mathbf A)^{-1}=(h_{ij})$, recoge los efectos directos e indirectos por toda la red \vertema{3.A.31}.
\item \emph{Perturbaciones de oferta: aguas abajo}. Ante cambios de productividad,
  \begin{equation*}
  \mathrm{d}\ln y_i=\underbrace{\mathrm{d}z_i}_{\text{efecto propio}}+\underbrace{\sum_{j}(h_{ij}-\mathbb 1_{j=i})\,\mathrm{d}z_j}_{\text{efecto red}}\quad\Longleftrightarrow\quad \mathrm{d}\ln\mathbf y=\mathbf H\,\mathrm{d}\mathbf z :
  \end{equation*}
  si cae la productividad de una industria, sube su precio y caen la producción de las que la usan como insumo.
\item \emph{Perturbaciones de demanda: aguas arriba}. Una caída del gasto en una industria reduce lo que compra a sus proveedores, y se propaga hacia ellos.
\item \emph{Un ejemplo con 4 industrias} (figura~\ref{figu:redes-produccion}): la 1 compra a la 2 y a la 4 y vende a la 3 y a la 4. Una perturbación de oferta en la 1 llega a la 3 y a la 4, no a la 2 (con \autor{Cobb}-\autor{Douglas}, el efecto precio y el efecto cantidad se compensan); una de demanda llega a la 2 y a la 4, no a la 3.
\item \emph{La evidencia}: con datos sectoriales de Estados Unidos de 1991 a 2009, el efecto de las perturbaciones a través de la red es mayor que el directo, un \emph{multiplicador de red}.
\end{itemize}

\begin{infografia}[H]
\centering
\caption{Una red de producción de 4 industrias}\label{figu:redes-produccion}
\grafico[0.55\textwidth]{redes-produccion}
\fuente{Elaboración propia a partir de \textcite{AcemogluAkcigitKerr2016}. Las flechas van del proveedor al comprador, con el peso $a_{ij}$ del insumo $j$ en la producción de $i$.}
\end{infografia}
\findemo

% ==========================================================================
\section{Anexo: perturbaciones y respuestas de política}\label{anexo:perturbaciones}

\begin{table}[H]
\centering
\caption{Perturbaciones, régimen cambiario y respuesta de política en una economía abierta}\label{tab:perturbaciones}
\small
\begin{tabular}{>{\raggedright\arraybackslash}m{3.2cm}>{\raggedright\arraybackslash}m{3.4cm}>{\raggedright\arraybackslash}m{3.7cm}>{\raggedright\arraybackslash}m{3.7cm}}
\toprule
\rowcolor{tceeazulpalido}\textbf{Perturbación} & \textbf{Efecto} & \textbf{Con tipo flexible} & \textbf{Con tipo fijo o unión monetaria}\\
\midrule
\textcolor{tceeazul}{\textbf{Caída de la demanda interna}} & Menos renta e inflación & Bajar los tipos con el tipo natural (divina coincidencia) & Política fiscal\\
\addlinespace[0.3em]
\textcolor{tceeazul}{\textbf{Caída de la demanda exterior}} & Menos exportaciones y renta & La depreciación la compensa (aislamiento) & Cae la renta: política fiscal\\
\addlinespace[0.3em]
\textcolor{tceeazul}{\textbf{Subida del tipo de interés mundial}} & Salida de capitales & Depreciación; por el canal financiero, contracción & Sube el tipo propio: cae la inversión\\
\addlinespace[0.3em]
\textcolor{tceerojo}{\textbf{Encarecimiento de la energía importada}} & Estanflación; inflación importada & Mirar a través del efecto directo y responder a la segunda ronda; la depreciación puede amplificarla & Política fiscal focalizada y temporal; política de oferta (energía)\\
\addlinespace[0.3em]
\textcolor{tceeverde}{\textbf{Mejora de la productividad}} & Más producción, menos precios & Política monetaria más expansiva: $\pi_{H}=0$ & El ajuste de los precios relativos, por los precios\\
\addlinespace[0.3em]
\textcolor{tceerojo}{\textbf{Pérdida de competitividad}} & Déficit exterior y paro & Depreciación nominal & Devaluación interna, devaluación fiscal, reformas\\
\bottomrule
\end{tabular}
\fuente{Elaboración propia a partir de \textcite{Mundell1963Capital}, \textcite{GaliMonacelli2005Monetary}, \textcite{BlanchardGali2010Oil} y \textcite{FarhiGopinathItskhoki2014}. En azul, perturbaciones de demanda; en rojo, de oferta adversas; en verde, de oferta favorables.}
\end{table}

% ==========================================================================
\section{Anexo: aspectos para el test}\label{anexo:test}

\begin{anotaciones}
Preguntas tipo test del primer ejercicio: \href{https://www.victorgutierrezmarcos.es/oposicion/temario/primer-ejercicio/test/simulador.html?tema=3.A.32}{simulador de test}.
\end{anotaciones}

\begin{itemize}
\item \emph{Fechas y obras}:
  \begin{itemize}
  \item \textcite{Fleming1962Domestic}, en las \emph{IMF Staff Papers}, y \textcite{Mundell1963Capital}, en el \emph{Canadian Journal of Economics and Political Science}: el modelo IS-LM-BP. \textcite{Mundell1962Appropriate}: la asignación de los instrumentos. \autor{Robert Mundell}, Nobel de 1999.
  \item \textcite{Meade1951}: equilibrio interno y externo (Nobel de 1977, con \autor{Bertil Ohlin}); \textcite{Tinbergen1952Theory}: tantos instrumentos como objetivos; \textcite{Swan1963Longer}, escrito en 1955.
  \item \textcite{Dornbusch1976Expectations}: sobrerreacción del tipo de cambio. \textcite{ObstfeldRogoff1995Redux}: nueva macroeconomía abierta. \textcite{ClaridaGaliGertler1999Science} y \textcite{Woodford2003Interest}: el modelo neokeynesiano (Nobel de 2026 a \autor{Woodford}, con \autor{Hall}). \textcite{GaliMonacelli2005Monetary}: economía pequeña neokeynesiana; \textcite{GaliMonacelli2008Optimal}: unión monetaria.
  \item \textcite{BlanchardQuah1989}: la oferta tiene efectos permanentes; \textcite{Gali1999Technology}: tecnología y horas; \textcite{Hamilton1983Oil}: petróleo y recesiones; \textcite{Kilian2009NotAll}: no todas las perturbaciones del petróleo son iguales; \textcite{BrunoSachs1985Economics}: estanflación y salarios reales rígidos; \textcite{GuerrieriEtAl2022}: perturbaciones de oferta keynesianas.
  \end{itemize}
\item \emph{Trampas típicas}:
  \begin{itemize}
  \item La BP vertical es movilidad \emph{nula} ($\beta_K=0$), no perfecta; su pendiente es $m/\beta_K$.
  \item Con tipo fijo, la política monetaria es ineficaz con \emph{cualquier} movilidad, porque el dinero es endógeno; lo que cambia es la rapidez con que se pierden las reservas.
  \item Con tipo flexible y movilidad perfecta, la fiscal es ineficaz porque expulsa \emph{exportaciones netas}, no inversión.
  \item Con tipo fijo, la DA abierta tiene pendiente negativa por el efecto competitividad; con flexible, por el efecto saldos reales.
  \item $E$ es el tipo directo: si sube, la moneda nacional se deprecia; $q=EP^*/P$ sube con la depreciación real.
  \item Con tipo flexible, la renta no depende de la demanda exterior en el modelo (aislamiento), pero sí del tipo de interés mundial.
  \item En el modelo neokeynesiano, una perturbación tecnológica desplaza la DA (por el tipo natural), no la NKPC, y no rompe la divina coincidencia; la rompen las de costes.
  \item En el modelo de \autor{Galí} y \autor{Monacelli}, la regla con la inflación interna domina a la de la inflación de consumo, y esta al tipo fijo; la divina coincidencia vale para $\pi_{H,t}$, no para la del consumo.
  \item Una perturbación de oferta mueve producción y precios en sentido contrario; una de demanda, en el mismo.
  \item El \emph{cuarteto imposible} (\autor{Padoa-Schioppa}, 1988) añade el libre comercio a los 3 objetivos del trilema.
  \item En \autor{Blanchard} y \autor{Quah} (1989), la demanda no tiene efectos a largo plazo sobre la producción; la oferta, sí.
  \end{itemize}
\end{itemize}

\bibliografia
\end{document}
